\documentclass[
  12pt,
  a4paper,
  oneside
]{scrreprt}

% =================================================
% Sprache, Schrift und Kodierung
% =================================================
\usepackage[T1]{fontenc}
\usepackage[utf8]{inputenc}
\usepackage[ngerman]{babel}
\usepackage{lmodern}
\usepackage{microtype}

% =================================================
% Seitenlayout
% =================================================
\usepackage[
  left=2.7cm,
  right=2.7cm,
  top=2.7cm,
  bottom=2.7cm
]{geometry}

\usepackage{setspace}
\onehalfspacing

% =================================================
% Mathematik, Tabellen und Listen
% =================================================
\usepackage{amsmath}
\usepackage{amssymb}
\usepackage{booktabs}
\usepackage{longtable}
\usepackage{array}
\usepackage{tabularx}
\usepackage{enumitem}
\usepackage{multirow}

% =================================================
% Gestaltung
% =================================================
\usepackage{xcolor}
\usepackage[most]{tcolorbox}
\usepackage{titlesec}
\usepackage{hyperref}

\hypersetup{
  colorlinks=true,
  linkcolor=black,
  urlcolor=blue,
  citecolor=black,
  pdftitle={Zusammenfassung Werkstoff- und Materialkunde}
}

\titleformat{\chapter}
  {\normalfont\Huge\bfseries}
  {\thechapter.}
  {0.7em}
  {}

\titleformat{\section}
  {\normalfont\Large\bfseries}
  {\thesection}
  {0.7em}
  {}

\setlist[itemize]{
  itemsep=0.35em,
  topsep=0.35em
}

\setlist[enumerate]{
  itemsep=0.35em,
  topsep=0.35em
}

% =================================================
% Eigene Befehle und Boxen
% =================================================
\newcommand{\quelle}[1]{\hfill\texttt{#1}}

\newtcolorbox{merksatzbox}{
  colback=gray!8,
  colframe=black!60,
  boxrule=0.6pt,
  arc=1.5mm,
  left=3mm,
  right=3mm,
  top=2mm,
  bottom=2mm
}

\newtcolorbox{darstellungsbox}{
  colback=blue!4,
  colframe=blue!50!black,
  boxrule=0.6pt,
  arc=1.5mm,
  left=3mm,
  right=3mm,
  top=2mm,
  bottom=2mm
}

% =================================================
% Beginn des Dokuments
% =================================================
\begin{document}

% =================================================
% Titelseite
% =================================================
\begin{titlepage}
  \centering
  \vspace*{3cm}

  {\Huge\bfseries Werkstoff- und Materialkunde\par}
  \vspace{0.5cm}
  {\Large Zusammenfassung\par}

  \vspace{1.5cm}

  \begin{tcolorbox}[
    colback=gray!8,
    colframe=black!60,
    width=0.88\textwidth,
    arc=1.5mm
  ]
  Diese Zusammenfassung basiert ausschließlich auf dem bereitgestellten
  Skriptum \textit{Werkstoffkunde 1 (1).pdf}. Die einzelnen Kapitel
  werden in derselben fachlichen Reihenfolge wie im Skriptum ergänzt.
  \end{tcolorbox}

  \vfill

  {\large LaTeX-Quelltext für eine PDF-Erstellung\par}
  \vspace{0.5cm}
  {\large \today\par}
\end{titlepage}

% =================================================
% Inhaltsverzeichnis
% =================================================
\tableofcontents
\clearpage

% =================================================
% KAPITEL 1: GRUNDLAGEN
% =================================================
\chapter{Grundlagen}

\section{Werkstoffwahl im Bauwesen}

Eine wesentliche Aufgabe von Ingenieurinnen und Ingenieuren besteht
darin, für unterschiedliche Aufgaben den passenden Werkstoff
auszuwählen. Voraussetzung für eine korrekte Auswahl ist die Kenntnis
der verfügbaren Baustoffe und ihrer Eigenschaften. 

Das Dokument behandelt unter anderem folgende Werk- und Baustoffe:

\begin{itemize}
  \item Holz und Holzwerkstoffe,
  \item Gesteinskörnung,
  \item Zement,
  \item Beton,
  \item metallische Werkstoffe,
  \item Glas,
  \item Asphalt,
  \item Mauerwerk,
  \item neue Materialien,
  \item biobasierte und kreislauffähige Baustoffe.
\end{itemize}



Die Wahl eines Baustoffs hängt nicht nur von seinen physikalischen und
mechanischen Eigenschaften ab. Auch ökonomische und ökologische
Kriterien müssen berücksichtigt werden. In der Vergangenheit war zudem
die Verfügbarkeit eines Baustoffs ein wesentliches Kriterium.


\begin{merksatzbox}
\textbf{Merksatz:}
Die Auswahl eines Baustoffs beruht auf physikalischen, mechanischen,
ökonomischen und ökologischen Kriterien. Zusätzlich ist die Kenntnis
der verfügbaren Werkstoffe und ihrer Eigenschaften erforderlich.

\end{merksatzbox}

\section{Einteilung der Baustoffe}

Baustoffe können grundsätzlich in anorganische und organische
Werkstoffe eingeteilt werden. 

\begin{longtable}{
  >{\raggedright\arraybackslash}p{0.38\textwidth}
  >{\raggedright\arraybackslash}p{0.38\textwidth}
}
\toprule
\textbf{Anorganische Baustoffe} &
\textbf{Organische Baustoffe} \\
\midrule
\endhead

Naturstein &
Holz \\

Glas &
Holzwerkstoffe \\

Keramische Baustoffe &
Bitumen \\

Mörtel &
Asphalt \\

Beton &
Faserverbundkunststoffe \\

Lehm &
Kunststoffe \\

Gusseisen, Stahl und Aluminium &
\\

\bottomrule
\end{longtable}



Im Bauwesen werden Werkstoffe häufig kombiniert, weil jeder Baustoff
unterschiedliche Vor- und Nachteile besitzt. Als Beispiel nennt das
Dokument Stahlbeton: Beton übernimmt vor allem die Abtragung von
Druckkräften, während die Betonstahlbewehrung die Zugkraftübertragung
übernimmt. 

\section{Kraft}

Kraft ist eine wesentliche physikalische Größe. Nach dem zweiten
Newton'schen Axiom ergibt sich die Kraft aus dem Produkt der Masse $m$
und der Beschleunigung $a$:

\begin{equation}
  F = m \cdot a
\end{equation}

Dabei bedeuten:

\begin{itemize}
  \item $F$: Kraft in $\mathrm{N}$,
  \item $m$: Masse in $\mathrm{kg}$,
  \item $a$: Beschleunigung in $\mathrm{m/s^2}$.
\end{itemize}



Neben Newton werden für größere Kräfte auch Kilonewton und Meganewton
verwendet:

\begin{align}
  1\,\mathrm{kN} &= 10^3\,\mathrm{N}, \\
  1\,\mathrm{MN} &= 10^6\,\mathrm{N}.
\end{align}



\section{Dichte}

Die Dichte beschreibt, wie viel Masse in einem bestimmten Volumen
enthalten ist. Sie wird mit dem Formelzeichen $\rho$ angegeben:

\begin{equation}
  \rho = \frac{m}{V}
\end{equation}

Dabei bedeuten:

\begin{itemize}
  \item $\rho$: Dichte in $\mathrm{kg/m^3}$,
  \item $m$: Masse in $\mathrm{kg}$,
  \item $V$: Volumen in $\mathrm{m^3}$.
\end{itemize}



Je nach Berücksichtigung von Hohlräumen und Poren wird zwischen
Reindichte, Rohdichte und Schüttdichte unterschieden.


\subsection{Reindichte}

Die Reindichte $\rho_0$ bezieht sich auf das Volumen eines Werkstoffs
ohne Hohlräume:

\begin{equation}
  \rho_0 = \frac{m}{V_0}
\end{equation}

Dabei ist $V_0$ das Volumen des Werkstoffs ohne Hohlräume.


\subsection{Rohdichte}

Die Rohdichte $\rho_r$ berücksichtigt zusätzlich das Volumen der Poren:

\begin{equation}
  \rho_r = \frac{m}{V_0 + V_p}
\end{equation}

Dabei ist $V_p$ das Volumen der Luft- und Kapillarporen.


\subsection{Schüttdichte}

Die Schüttdichte $\rho_s$ berücksichtigt zusätzlich die Zwischenräume
zwischen einzelnen Körnern eines Haufwerks:

\begin{equation}
  \rho_s = \frac{m}{V_0 + V_p + V_z}
\end{equation}

Dabei ist $V_z$ das Volumen der Zwischenräume.


\begin{darstellungsbox}
\textbf{Darstellung: Reindichte, Rohdichte und Schüttdichte}

Die Darstellung unterscheidet drei Bezugsvolumina:

\begin{itemize}
  \item Das Volumen eines vollständig porenfrei zerkleinerten Stoffes
  bildet die Grundlage für die Reindichte.

  \item Das Volumen eines einzelnen Korns einschließlich der Kornporen
  bildet die Grundlage für die Rohdichte.

  \item Das Volumen eines Prüfgefäßes mit einzelnen Körnern und
  Haufwerksporen bildet die Grundlage für die Schüttdichte.
\end{itemize}
\end{darstellungsbox}

\section{Zug und Druck}

Eine Zugkraft erzeugt Zugspannungen. Eine Druckkraft erzeugt
Druckspannungen. Die Zugspannung wird wie folgt berechnet:

\begin{equation}
  \sigma_Z = \frac{F}{A}
\end{equation}

Die Druckspannung wird mit negativem Vorzeichen angegeben:

\begin{equation}
  \sigma_D = -\frac{F}{A}
\end{equation}

Dabei bedeuten:

\begin{itemize}
  \item $\sigma_Z$: Zugspannung in $\mathrm{N/mm^2}$,
  \item $\sigma_D$: Druckspannung in $\mathrm{N/mm^2}$,
  \item $F$: Zug- oder Druckkraft in $\mathrm{N}$,
  \item $A$: belastete Querschnittsfläche in $\mathrm{mm^2}$.
\end{itemize}



Zugbeanspruchung führt zu einer Dehnung, Druckbeanspruchung zu einer
Stauchung. Zusätzlich tritt normal zur Belastungsrichtung eine
Querkontraktion auf. 

Die Dehnung wird mit folgender Gleichung beschrieben:

\begin{equation}
  \varepsilon = \frac{\Delta l}{l_0}
\end{equation}

Dabei bedeuten:

\begin{itemize}
  \item $\varepsilon$: Dehnung,
  \item $\Delta l$: Längenänderung,
  \item $l_0$: Ursprungslänge.
\end{itemize}



\section{Poisson'sche Zahl}

Das Verhältnis der Querkontraktion zur Längenänderung ist bei
elastischer Beanspruchung ein konstanter Werkstoffkennwert. Dieser
Kennwert wird Poisson'sche Zahl $\mu$ genannt:

\begin{equation}
  \mu = \frac{\varepsilon_q}{\varepsilon}
\end{equation}

Die Poisson'sche Zahl liegt laut Dokument zwischen $0$ und $0{,}5$.
Für Stahl wird ein Wert von $0{,}3$, für Beton ein Wert von $0{,}16$
angegeben. 

\section{Schubbeanspruchung}

Scher- beziehungsweise Schubkräfte sind gegenparallel wirkende Kräfte,
die auf eine parallele Fläche wirken. Sie erzeugen Schubspannungen.
Die dadurch entstehende Verformung wird als Schiebung oder Scherung
bezeichnet. 

Die Schubspannung $\tau$ ergibt sich aus:

\begin{equation}
  \tau = \frac{F}{A}
\end{equation}

Dabei bedeuten:

\begin{itemize}
  \item $\tau$: Schubspannung in $\mathrm{N/mm^2}$,
  \item $F$: Schubkraft in $\mathrm{N}$,
  \item $A$: Scherfläche in $\mathrm{mm^2}$.
\end{itemize}

Die Schiebung beziehungsweise Scherung kann beschrieben werden durch:

\begin{equation}
  \gamma = \frac{w}{l} = \tan(\vartheta)
\end{equation}

Dabei bedeuten:

\begin{itemize}
  \item $\gamma$: Schiebung beziehungsweise Scherung,
  \item $w$: Verschiebung in Richtung der Scherkraft,
  \item $l$: Höhe des Körpers,
  \item $\vartheta$: Scherwinkel.
\end{itemize}



\section{Biegung}

Ein auf Biegung beanspruchter Werkstoff muss ein Biegemoment aufnehmen.
Dieses verursacht Biegespannungen. Die Ermittlung der Biegespannungen
und Biegeverformungen hängt von der Belastungsart sowie vom statischen
System ab. 

Das Dokument unterscheidet insbesondere den Drei-Punkt- und den
Vier-Punkt-Biegeversuch:

\begin{itemize}
  \item Beim \textbf{Drei-Punkt-Biegeversuch} wirkt eine Einzellast in
  der Mitte des Prüfkörpers. Das maximale Moment und die maximale
  Spannung treten an einem Punkt auf.

  \item Beim \textbf{Vier-Punkt-Biegeversuch} wirken zwei Lasten im
  Feld. Zwischen diesen beiden Lasten ist das Moment konstant.
\end{itemize}



\begin{darstellungsbox}
\textbf{Darstellung: Drei-Punkt- und Vier-Punkt-Biegeversuch}

Beim Drei-Punkt-Biegeversuch liegt ein Prüfkörper auf zwei Lagern und
wird in der Mitte durch eine einzelne Last belastet.

Beim Vier-Punkt-Biegeversuch liegt der Prüfkörper ebenfalls auf zwei
Lagern. Zusätzlich wirken zwei Lasten im Feld. Dadurch entsteht
zwischen den beiden Lasten ein Bereich mit konstantem Biegemoment.
\end{darstellungsbox}

\section{Ökologische Eigenschaften von Baustoffen}

Das Dokument behandelt ökologische Eigenschaften von Baustoffen unter
anderem anhand der Lebenszyklusanalyse sowie anhand ausgewählter
Umwelteinwirkungen. 

Für Holz wird ein negatives Treibhausgaspotenzial in der Herstellung
angeführt. Dies wird auf die Kohlendioxidspeicherfähigkeit des Holzes
zurückgeführt. Bei Betrachtung des gesamten Lebenszyklus kann sich
diese Bewertung jedoch relativieren. 

Der Einsatz von Klebstoffen führt bei keilgezinktem Holz und
Brettschichtholz zu höheren Werten für Ozonabbaupotenzial und
nicht-erneuerbaren Primärenergiebedarf als bei Vollholz ohne
Keilzinkung. 

Mauerziegel besitzt laut Dokument ein vergleichbares
Treibhausgaspotenzial und Ozonabbaupotenzial wie Beton der
Festigkeitsklasse C20/25. Der nicht-erneuerbare Primärenergiebedarf
von Mauerziegeln wird jedoch als doppelt so hoch angegeben.


\begin{merksatzbox}
\textbf{Merksatz:}
Die Umwelteinwirkungen von Baustoffen sind nicht allein anhand von
Masse oder Volumen vergleichbar. Laut Dokument kann die Einbeziehung
von Werkstoffeigenschaften wie Festigkeiten die Bewertung der
Umwelteinwirkungen relativieren. 
\end{merksatzbox}

% =================================================
% HIER WERDEN DIE NÄCHSTEN KAPITEL EINGEFÜGT:
%
% Kapitel 2: Beton
% Kapitel 3: Bindemittel und Zement
% Kapitel 4: Gestein, Naturwerkstein und Gesteinskörnung
% Kapitel 5: Metallische Werkstoffe
% Kapitel 6: Glas
% Kapitel 7: Holz und Holzwerkstoffe
% Kapitel 8: Asphalt und Straßenbau
% Kapitel 9: Mauerwerk
% Kapitel 10: Neue, biobasierte und kreislauffähige Materialien
% =================================================
% =================================================
% KAPITEL 2: BETON
% Diesen gesamten Block vor \end{document} einfügen.
% =================================================

\chapter{Beton}

\section{Definition und Zusammensetzung}

Beton ist ein erhärtetes Gemisch aus Zement, grober und feiner
Gesteinskörnung sowie Wasser. Zusatzmittel und Zusatzstoffe können
dem Beton zusätzlich beigegeben werden. Beton zählt weltweit neben
Lehm zu den am häufigsten verwendeten Baustoffen.


Die im Dokument behandelten Bestandteile des Betons sind:

\begin{itemize}
  \item Zement,
  \item Gesteinskörnung,
  \item Zugabewasser,
  \item Luftporen,
  \item Zusatzstoffe,
  \item Zusatzmittel.
\end{itemize}



\begin{merksatzbox}
\textbf{Merksatz:}
Beton ist ein erhärtetes Gemisch aus Zement, grober und feiner
Gesteinskörnung sowie Wasser. Zusatzmittel und Zusatzstoffe können
zusätzlich enthalten sein.

\end{merksatzbox}

\section{Wasser-Bindemittel-Wert}

Der Wasser-Bindemittel-Wert, kurz W/B-Wert, beschreibt das Verhältnis
zwischen der Masse des Wassers und der Masse des Bindemittels.

\begin{equation}
  \text{W/B-Wert}
  =
  \frac{\text{Masse Wasser}}{\text{Masse Bindemittel}}
\end{equation}



Für die vollständige Hydratation des Zements wird ein W/B-Wert von
$0{,}4$ angegeben. Bei einem W/B-Wert von mehr als $0{,}4$ bleibt
Wasser in Kapillarporen zurück. Dies führt zu einer geringeren
Festigkeit, Beständigkeit, Dichtheit und einem geringeren
Korrosionsschutz.


\begin{darstellungsbox}
\textbf{Darstellung: Einfluss des W/B-Werts}

Mit zunehmendem Wasser-Bindemittel-Wert nimmt die Betonfestigkeit ab.
Gleichzeitig nimmt der Anteil an Kapillarporen zu. Ein W/B-Wert von
mehr als $0{,}4$ führt laut Dokument dazu, dass Wasser in
Kapillarporen zurückbleiben kann.
\end{darstellungsbox}

\section{Frischbeton}

Frischbeton ist Beton vor dem Erhärten. Zu den im Dokument behandelten
Eigenschaften des Frischbetons zählen:

\begin{itemize}
  \item Frischbetonrohdichte,
  \item Konsistenz,
  \item Luftporengehalt,
  \item Wärmeentwicklung beim Abbinden.
\end{itemize}



\subsection{Frischbetonrohdichte}

Zur Bestimmung der Frischbetonrohdichte wird Frischbeton in einem
biegesteifen und wasserdichten Behälter verdichtet und gewogen. Die
Rohdichte wird aus Masse und Volumen ermittelt:

\begin{equation}
  \rho = \frac{m}{V}
\end{equation}

Dabei bedeuten:

\begin{itemize}
  \item $\rho$: Rohdichte,
  \item $m$: Masse,
  \item $V$: Volumen.
\end{itemize}

Die Rohdichte hängt hauptsächlich von der Gesteinskörnung und zusätzlich
vom Luftporengehalt ab. Für Beton wird eine durchschnittliche Dichte
von ungefähr

\begin{equation}
  2400\,\mathrm{kg/m^3}
\end{equation}

angegeben.


\subsection{Konsistenz}

Die Konsistenz beschreibt die Verarbeitbarkeit von Frischbeton. Für
steifere Konsistenzklassen wird das Verdichtungsmaß bestimmt.

\begin{equation}
  C = \frac{h_1}{h}
\end{equation}

Dabei bedeuten:

\begin{itemize}
  \item $C$: Verdichtungsmaß,
  \item $h_1$: Höhe des lose eingefüllten Betons,
  \item $h$: Höhe des verdichteten Betons.
\end{itemize}



Für dünnflüssigere Konsistenzklassen wird das Ausbreitmaß verwendet.
Dabei wird Frischbeton auf einem Ausbreittisch geprüft. Nach dem
Anheben und Fallenlassen des Tisches werden zwei Durchmesser des
Betonkuchens gemessen.

\begin{equation}
  F = \frac{d_1 + d_2}{2}
\end{equation}

Dabei bedeuten:

\begin{itemize}
  \item $F$: Ausbreitmaß,
  \item $d_1$: erster gemessener Durchmesser,
  \item $d_2$: zweiter gemessener Durchmesser.
\end{itemize}



\subsection{Luftporengehalt}

Der Luftporengehalt von Frischbeton wird mit einem Luftporentopf
ermittelt. Als übliche Luftporengehalte werden Werte von

\begin{equation}
  1 \text{ bis } 2\,\mathrm{Vol.\%}
\end{equation}

angegeben.


\subsection{Wärmeentwicklung beim Abbinden}

Beim Abbinden von Beton entsteht Wärme. Diese Wärmeentwicklung ist
insbesondere bei massigen Bauteilen von Bedeutung, weil
Temperaturrisse vermieden werden sollen.


Die maximale Betontemperatur im Anlieferzustand beträgt laut Dokument:

\begin{equation}
  27\,^\circ\mathrm{C}
\end{equation}

Für Beton der Weißen Wanne wird eine maximale Temperatur von

\begin{equation}
  22\,^\circ\mathrm{C}
\end{equation}

angegeben.


\section{Festbeton}

Nach dem Erhärten wird Beton als Festbeton bezeichnet. Die
Eigenschaften des Festbetons umfassen laut Dokument insbesondere:

\begin{itemize}
  \item Druckfestigkeit,
  \item Zugfestigkeit,
  \item elastische Verformung beziehungsweise E-Modul.
\end{itemize}



\subsection{Druckfestigkeit}

Beton kann vor allem Druckspannungen aufnehmen. Die Einteilung von
Beton erfolgt unter anderem anhand von Druckfestigkeitsklassen. Die
Druckfestigkeit wird an Prüfkörpern bestimmt und bezieht sich auf ein
Betonalter von 28 Tagen.


Als Beispiel für eine Betonbezeichnung wird angeführt:

\begin{center}
\texttt{C25/30 XC3/XD2/XF3/XA1L/SB}
\end{center}

Dabei bezeichnet \texttt{C25/30} die Druckfestigkeitsklasse. Die
weiteren Kürzel bezeichnen Expositionsklassen. \texttt{SB} wird als
Betonanwendung Sichtbeton angeführt.


\subsection{Elastizitätsmodul von Beton}

Der Elastizitätsmodul beschreibt den Zusammenhang zwischen Spannung und
Dehnung. Er ist ein Maß für den Widerstand eines Werkstoffs gegen
elastische Verformung.


Der statische Elastizitätsmodul des Betons wird mit folgender Beziehung
ermittelt:

\begin{equation}
  E_b =
  \frac{\sigma_a - \sigma_b}
  {\varepsilon_a - \varepsilon_b}
\end{equation}

Dabei bedeuten:

\begin{itemize}
  \item $E_b$: statischer Elastizitätsmodul,
  \item $\sigma_a$: obere Spannung,
  \item $\sigma_b$: untere Spannung,
  \item $\varepsilon_a$: mittlere Dehnung unter $\sigma_a$,
  \item $\varepsilon_b$: mittlere Dehnung unter $\sigma_b$.
\end{itemize}



Für Konstruktionsbeton soll der E-Modul groß sein, damit die
Verformungen klein bleiben. Für Betonstraßen und Massenbeton soll der
E-Modul eher niedrig sein, um hohe Spannungen bei auftretenden
Dehnungen zu vermeiden.


\begin{darstellungsbox}
\textbf{Darstellung: Prüfung des E-Moduls}

Zur Bestimmung des E-Moduls wird ein Prüfkörper in eine Druckpresse
eingespannt. Verformungsmessgeräte erfassen die auftretenden Dehnungen.
Anschließend wird ein Belastungszyklus durchgeführt. Der E-Modul wird
aus dem Spannungs-Dehnungs-Verhalten berechnet.
\end{darstellungsbox}

\section{Expositionsklassen}

Expositionsklassen sind ein zentrales Werkzeug der EN~206-1. Sie
definieren Angriffsrisiken durch Umwelteinflüsse. Die
Mindestanforderungen in Abhängigkeit vom Risiko werden in
ÖNORM~B~4710-1 beschrieben.


\begin{longtable}{
  >{\centering\arraybackslash}p{0.18\textwidth}
  >{\raggedright\arraybackslash}p{0.56\textwidth}
}
\toprule
\textbf{Klasse} & \textbf{Bedeutung} \\
\midrule
\endhead

X0 & kein Angriffsrisiko \\

XC & Korrosion durch Karbonatisierung \\

XD & Chloride, beispielsweise durch Taumittel \\

XS & Meerwasser \\

XF & Frost und Frost-Tausalz \\

XA & chemischer Angriff \\

XM & Verschleiß beziehungsweise mechanischer Abrieb \\

\bottomrule
\end{longtable}



\subsection{Karbonatisierung}

Beton besitzt einen pH-Wert von ungefähr 12. Dadurch wird
Bewehrungsstahl durch eine Passivierungsschicht geschützt. Durch
Karbonatisierung sinkt der pH-Wert. Dadurch kann die Bewehrung zu
rosten beginnen.


Die im Dokument angeführte Reaktion lautet:

\begin{equation}
  \mathrm{
  Ca(OH)_2 + H_2O + CO_2
  \rightarrow
  CaCO_3 + H_2O
  }
\end{equation}

Die Karbonatisierungstiefe kann mit dem pH-Indikator Phenolphthalein
überprüft werden.


\section{Stoffraumrechnung}

Die Stoffraumrechnung dient der rechnerischen Bestimmung der
Zusammensetzung von $1\,\mathrm{m^3}$ Frischbeton.

\begin{equation}
  1000 =
  \frac{z}{\rho_z}
  +
  \frac{f}{\rho_f}
  +
  \frac{w}{\rho_w}
  +
  \frac{g}{\rho_g}
  +
  p
  \qquad [\mathrm{dm^3/m^3}]
\end{equation}

Dabei bedeuten:

\begin{longtable}{
  >{\centering\arraybackslash}p{0.15\textwidth}
  >{\raggedright\arraybackslash}p{0.56\textwidth}
}
\toprule
\textbf{Symbol} & \textbf{Bedeutung} \\
\midrule
\endhead

$z$ & Zementgehalt in $\mathrm{kg/m^3}$ \\

$\rho_z$ & Dichte des Zements in $\mathrm{kg/dm^3}$ \\

$f$ & Zusatzstoffgehalt in $\mathrm{kg/m^3}$ \\

$\rho_f$ & Dichte des Zusatzstoffs in $\mathrm{kg/dm^3}$ \\

$w$ & Wassergehalt in $\mathrm{kg/m^3}$ \\

$\rho_w$ & Dichte des Wassers in $\mathrm{kg/dm^3}$ \\

$g$ & Gehalt der Gesteinskörnung in $\mathrm{kg/m^3}$ \\

$\rho_g$ & Rohdichte der Gesteinskörnung in $\mathrm{kg/dm^3}$ \\

$p$ & Porenvolumen in $\mathrm{dm^3/m^3}$ \\

\bottomrule
\end{longtable}



\section{Betonherstellung}

Beton wird großtechnisch in Betonmischanlagen hergestellt. Als
Bestandteile einer Betonmischanlage werden genannt:

\begin{itemize}
  \item Lagereinrichtung für Gesteinskörnung,
  \item Zementsilo,
  \item Zugabeeinrichtung für Betonzusatzmittel,
  \item Steuerungstechnik,
  \item Dosiereinrichtungen,
  \item Betonmischer,
  \item Übergabeeinrichtung für Frischbeton an einen LKW.
\end{itemize}



\section{Einbau, Verdichtung und Nachbehandlung}

Die freie Fallhöhe beim Betonieren darf maximal

\begin{equation}
  1{,}5\,\mathrm{m}
\end{equation}

betragen, da es sonst zu einer Entmischung kommen kann.


Als Verdichtungsverfahren nennt das Dokument:

\begin{itemize}
  \item Innenrüttler,
  \item Schalungsrüttler,
  \item Oberflächenrüttler wie Rüttelbohlen und Rüttelplatten.
\end{itemize}



Ab der Konsistenzklasse F66 ist laut Dokument keine Verdichtung mehr
erforderlich.


Die Nachbehandlung schützt Beton vor vorzeitigem Austrocknen, starken
Temperaturveränderungen und Rissbildung. Als Maßnahmen werden Folien
und Matten genannt.


\begin{merksatzbox}
\textbf{Merksatz:}
Verdichtung und Nachbehandlung sind wesentliche Schritte für die
Betonqualität. Die Nachbehandlung schützt insbesondere vor vorzeitigem
Austrocknen, starken Temperaturänderungen und Rissbildung.

\end{merksatzbox}
% =================================================
% KAPITEL 3: BINDEMITTEL UND ZEMENT
% Diesen gesamten Block vor \end{document} einfügen.
% =================================================

\chapter{Bindemittel und Zement}

\section{Bindemittel im Beton}

Bindemittel verbinden die einzelnen Bestandteile eines Baustoffs.
Im Beton übernimmt Zement die Funktion des Bindemittels. Beton besteht
aus Zement, Gesteinskörnung, Wasser sowie gegebenenfalls Zusatzmitteln
und Zusatzstoffen. 

Der Wasser-Bindemittel-Wert beschreibt das Verhältnis von Wasser zu
Bindemittel:

\begin{equation}
  \text{W/B-Wert}
  =
  \frac{\text{Masse Wasser}}
  {\text{Masse Bindemittel}}
\end{equation}

Zusatzstoffe dürfen laut Dokument nicht immer vollständig auf den
Bindemittelgehalt angerechnet werden. Dafür wird ein $k$-Faktor
verwendet:

\begin{equation}
  B =
  \text{Zement}
  +
  k \cdot \text{Zusatzstoff}
\end{equation}

Dabei ist $B$ der anrechenbare Bindemittelgehalt.


\begin{longtable}{
  >{\raggedright\arraybackslash}p{0.38\textwidth}
  >{\centering\arraybackslash}p{0.24\textwidth}
}
\toprule
\textbf{Zusatzstoff} & \textbf{$k$-Faktor} \\
\midrule
\endhead

Flugasche & $0{,}8$ \\

Hüttensand & $0{,}8$ \\

Silicastaub & $2{,}0$ \\

\bottomrule
\end{longtable}



\begin{merksatzbox}
\textbf{Merksatz:}
Der anrechenbare Bindemittelgehalt setzt sich aus Zement und dem mit
einem $k$-Faktor gewichteten Zusatzstoff zusammen. Für Flugasche und
Hüttensand gilt laut Dokument $k=0{,}8$, für Silicastaub $k=2{,}0$.

\end{merksatzbox}

\section{Normfestigkeitsklassen von Zement}

Die Normfestigkeitsklassen von Zement beziehen sich auf die
Druckfestigkeit nach 28 Tagen. Das Dokument nennt die
Normfestigkeitsklassen:

\begin{itemize}
  \item $32{,}5\,\mathrm{N/mm^2}$,
  \item $42{,}5\,\mathrm{N/mm^2}$,
  \item $52{,}5\,\mathrm{N/mm^2}$.
\end{itemize}



Zusätzlich wird die Anfangsfestigkeit unterschieden:

\begin{itemize}
  \item \texttt{N}: normale Anfangsfestigkeit,
  \item \texttt{R}: hohe beziehungsweise rapide Anfangsfestigkeit.
\end{itemize}



\begin{longtable}{
  >{\centering\arraybackslash}p{0.14\textwidth}
  >{\centering\arraybackslash}p{0.24\textwidth}
  >{\centering\arraybackslash}p{0.22\textwidth}
  >{\centering\arraybackslash}p{0.19\textwidth}
}
\toprule
\textbf{Festigkeitsklasse} &
\textbf{Anfangsfestigkeit} &
\textbf{Normfestigkeit nach 28 Tagen} &
\textbf{Erstarrungsbeginn} \\
\midrule
\endhead

32,5 N &
mindestens $16{,}0\,\mathrm{N/mm^2}$ nach 7 Tagen &
mindestens $32{,}5\,\mathrm{N/mm^2}$ &
mindestens $75\,\mathrm{min}$ \\

32,5 R &
mindestens $10{,}0\,\mathrm{N/mm^2}$ nach 2 Tagen &
mindestens $32{,}5\,\mathrm{N/mm^2}$ &
mindestens $75\,\mathrm{min}$ \\

42,5 N &
mindestens $10{,}0\,\mathrm{N/mm^2}$ nach 2 Tagen &
mindestens $42{,}5\,\mathrm{N/mm^2}$ &
mindestens $60\,\mathrm{min}$ \\

42,5 R &
mindestens $20{,}0\,\mathrm{N/mm^2}$ nach 2 Tagen &
mindestens $42{,}5\,\mathrm{N/mm^2}$ &
mindestens $60\,\mathrm{min}$ \\

52,5 N &
mindestens $20{,}0\,\mathrm{N/mm^2}$ nach 2 Tagen &
mindestens $52{,}5\,\mathrm{N/mm^2}$ &
mindestens $45\,\mathrm{min}$ \\

52,5 R &
mindestens $30{,}0\,\mathrm{N/mm^2}$ nach 2 Tagen &
mindestens $52{,}5\,\mathrm{N/mm^2}$ &
mindestens $45\,\mathrm{min}$ \\

\bottomrule
\end{longtable}

Die Tabelle beschreibt Anfangs- und Normfestigkeiten von Zement nach
ÖNORM EN 197-1. 

\begin{merksatzbox}
\textbf{Merksatz:}
Die Normfestigkeit von Zement bezieht sich auf die Druckfestigkeit nach
28 Tagen. Das Kürzel \texttt{N} steht für normale Anfangsfestigkeit,
während \texttt{R} eine hohe beziehungsweise rapide Anfangsfestigkeit
bezeichnet. 
\end{merksatzbox}

\section{Mahlfeinheit und Erhärtungsgeschwindigkeit}

Die Erhärtungsgeschwindigkeit von Zement hängt laut Dokument
insbesondere von der Mahlfeinheit ab. Je feiner ein Zement gemahlen
ist, desto größer ist seine Oberfläche und desto reaktiver ist er.


Die spezifische Oberfläche wird mit dem Blaine-Wert in
$\mathrm{cm^2/g}$ beschrieben.

\begin{longtable}{
  >{\raggedright\arraybackslash}p{0.44\textwidth}
  >{\centering\arraybackslash}p{0.30\textwidth}
}
\toprule
\textbf{Material} &
\textbf{Spezifische Oberfläche} \\
\midrule
\endhead

Standardzement &
ca.\ $3500\,\mathrm{cm^2/g}$ \\

Frühhochfester Zement &
ca.\ $5000\,\mathrm{cm^2/g}$ \\

Mikrosilika &
ca.\ $200000\,\mathrm{cm^2/g}$ \\

\bottomrule
\end{longtable}



\begin{darstellungsbox}
\textbf{Darstellung: Festigkeitsentwicklung von Zement}

Die Darstellung unterscheidet die Anfangserhärtung sowie die Erhärtung
und Nacherhärtung. Sie zeigt einen Verlauf für frühfesten Beton und
einen Verlauf für normalerhärtenden Beton.

Frühfester Beton erreicht in der frühen Phase schneller eine höhere
Festigkeit. Die 28-Tage-Festigkeit bleibt jedoch der maßgebende Wert
für die Zuordnung zur Normfestigkeitsklasse.
\end{darstellungsbox}

\section{Zementbezeichnung}

Die vollständige Bezeichnung eines Zements enthält laut Dokument:

\begin{itemize}
  \item die Norm,
  \item den Zementtyp beziehungsweise die Zusammensetzung,
  \item die Druckfestigkeitsklasse,
  \item die Frühfestigkeit.
\end{itemize}



Als Beispiel wird folgende vollständige Zementbezeichnung angegeben:

\begin{center}
\texttt{EN 197-1 -- CEM II/A-S 32,5 R}
\end{center}

\begin{longtable}{
  >{\raggedright\arraybackslash}p{0.28\textwidth}
  >{\raggedright\arraybackslash}p{0.52\textwidth}
}
\toprule
\textbf{Bestandteil} &
\textbf{Bedeutung} \\
\midrule
\endhead

\texttt{EN 197-1} &
Norm \\

\texttt{CEM II/A-S} &
Zementtyp; Kennzeichnung für Zusammensetzung und Zumahlstoffe \\

\texttt{32,5} &
Druckfestigkeitsklasse \\

\texttt{R} &
hohe beziehungsweise rapide Anfangsfestigkeit \\

\bottomrule
\end{longtable}



Für Zumahlstoffe werden im Dokument folgende Kennbuchstaben genannt:

\begin{longtable}{
  >{\centering\arraybackslash}p{0.20\textwidth}
  >{\raggedright\arraybackslash}p{0.50\textwidth}
}
\toprule
\textbf{Kennbuchstabe} &
\textbf{Zumahlstoff} \\
\midrule
\endhead

S & Hüttensand \\

V & Flugasche \\

L & Kalkstein \\

M & Kombination von Zumahlstoffen \\

\bottomrule
\end{longtable}



\section{Normenmörtel}

Zur Prüfung von Zement wird ein Normenmörtel hergestellt. Die im
Dokument angegebene Zusammensetzung lautet:

\begin{longtable}{
  >{\raggedright\arraybackslash}p{0.42\textwidth}
  >{\centering\arraybackslash}p{0.26\textwidth}
}
\toprule
\textbf{Bestandteil} &
\textbf{Masse} \\
\midrule
\endhead

Wasser &
$225\,\mathrm{g}$ \\

Zement &
$450\,\mathrm{g}$ \\

Normensand &
$1350\,\mathrm{g}$ \\

\bottomrule
\end{longtable}



Der Normenmörtel wird in Formen mit folgenden Abmessungen eingefüllt:

\begin{equation}
  40 \times 40 \times 160\,\mathrm{mm}
\end{equation}

Die Prüfkörper werden verdichtet und abgestrichen. 

\section{Lagerung der Zementprüfkörper}

Die Lagerung der Zementprüfkörper erfolgt laut Dokument in zwei
Schritten:

\begin{enumerate}
  \item Zunächst werden die Prüfkörper 24 Stunden lang in einem
  Feuchtschrank bei $20\,^\circ\mathrm{C}$ und $90\,\%$ relativer
  Feuchte gelagert.

  \item Anschließend werden die Prüfkörper bis zum Alter von 28 Tagen
  in einem Wasserbecken gelagert.
\end{enumerate}



\section{Festigkeitsprüfung von Zement}

Nach 28 Tagen erfolgt die Festigkeitsprüfung des Normenmörtels. Das
Dokument nennt folgende Prüfungen:

\begin{itemize}
  \item Biegezugfestigkeitsprüfung im Drei- oder
  Vier-Punkt-Biegeversuch,
  \item Druckfestigkeitsprüfung im Druckversuch.
\end{itemize}



\begin{darstellungsbox}
\textbf{Darstellung: Ablauf der Zementprüfung}

Die Zementprüfung umfasst folgende Schritte:

\begin{enumerate}
  \item Herstellung des Normenmörtels aus Wasser, Zement und
  Normensand;
  \item Einfüllen in Formen mit den Abmessungen
  $40 \times 40 \times 160\,\mathrm{mm}$;
  \item Verdichten und Abstreichen;
  \item Lagerung für 24 Stunden im Feuchtschrank bei
  $20\,^\circ\mathrm{C}$ und $90\,\%$ relativer Feuchte;
  \item Lagerung im Wasserbecken bis zum Alter von 28 Tagen;
  \item Biegezug- und Druckfestigkeitsprüfung.
\end{enumerate}
\end{darstellungsbox}

\section{Zusammenfassung}

Zement ist das Bindemittel im Beton. Seine maßgeblichen Eigenschaften
werden unter anderem durch die Normfestigkeitsklasse, die
Anfangsfestigkeit, die Mahlfeinheit und die Zusammensetzung bestimmt.
Die Zementprüfung erfolgt anhand eines genormten Mörtels und bewertet
insbesondere Biegezug- und Druckfestigkeit nach 28 Tagen.


\begin{merksatzbox}
\textbf{Merksatz:}
Für die Zementprüfung wird ein Normenmörtel aus $225\,\mathrm{g}$
Wasser, $450\,\mathrm{g}$ Zement und $1350\,\mathrm{g}$ Normensand
hergestellt. Die Prüfkörper werden bis zum Alter von 28 Tagen gelagert
und anschließend auf Biegezug- und Druckfestigkeit geprüft.

\end{merksatzbox}
% =================================================
% KAPITEL 4: GESTEIN, NATURWERKSTEIN
% UND GESTEINSKÖRNUNG
% Diesen gesamten Block vor \end{document} einfügen.
% =================================================

\chapter{Gestein, Naturwerkstein und Gesteinskörnung}

\section{Bedeutung von Gestein im Bauwesen}

Gestein und Gesteinskörnung sind bedeutende Baustoffe im Hochbau,
Tiefbau und Straßenbau. Sie werden unter anderem als Bestandteil von
Beton, Asphalt, ungebundenen Tragschichten und gebundenen
Tragschichten eingesetzt. 

Gesteinskörnungen werden außerdem mit Bindemitteln verkittet und für
Beton oder Asphalt verwendet. Die Einhaltung normativer Anforderungen
dient dazu, hochwertige und dauerhafte Gesteinskörnung bereitzustellen.


Im Straßenbau werden Gesteinskörnungen beispielsweise für Asphalt,
Beton, ungebundene Tragschichten und Frostschutzschichten benötigt.
Beim Bau der Umfahrung Wieselburg B25 wurden laut Dokument unter
anderem Gesteinskörnungen, Bindemittel und Beton eingesetzt.


\begin{merksatzbox}
\textbf{Merksatz:}
Gesteinskörnung ist ein wesentlicher Bestandteil von Beton, Asphalt,
gebundenen Tragschichten und ungebundenen Tragschichten. Ihre
Eigenschaften beeinflussen die Dauerhaftigkeit und die
Gebrauchseigenschaften des jeweiligen Baustoffs.

\end{merksatzbox}

\section{Naturwerkstein}

Naturwerkstein wird als natürlicher Baustoff eingesetzt. Im
bereitgestellten Dokument wird Gestein unter anderem im Zusammenhang
mit Naturwerkstein und Gesteinskörnung behandelt. Naturwerkstein kann
im Bauwesen beispielsweise für Bauwerke, Stützkonstruktionen,
Pflasterungen oder andere Anwendungen aus massivem Gestein verwendet
werden.

Für Bauvorhaben sind die Eigenschaften des anstehenden Gesteins von
besonderer Bedeutung. Im Zusammenhang mit dem Beispiel der Umfahrung
Wieselburg B25 werden unter anderem Einschnitte in Festgestein,
Wasserverhältnisse, die Fundierung einer Wildbrücke und die Fundierung
von Hochspannungsmasten genannt. 

Die Eignung eines Gesteins für eine Anwendung hängt insbesondere von
seiner Dauerhaftigkeit, seiner Korn- beziehungsweise Gefügeeigenschaft
und seiner Widerstandsfähigkeit gegen Beanspruchungen ab.


\section{Gesteinskörnung}

Gesteinskörnung kann aus natürlichen, industriell hergestellten oder
rezyklierten Materialien bestehen. Für ungebundene Tragschichten werden
Gemische aus natürlichen, industriell hergestellten und rezyklierten
Gesteinskörnungen verwendet. 

Gesteinskörnungen müssen unter anderem folgende Anforderungen erfüllen:

\begin{itemize}
  \item ausreichender Haftverbund mit Zementstein beziehungsweise
  Bitumen;
  \item keine schädlichen Verbindungen mit Zement, beispielsweise
  durch Sulfate;
  \item Gewährleistung der Beständigkeit von Beton oder Asphalt;
  \item keine Beeinträchtigung des Korrosionsschutzes der
  Betonbewehrung.
\end{itemize}



\section{Geometrische Anforderungen}

Die geometrischen Anforderungen an Gesteinskörnung betreffen
insbesondere:

\begin{itemize}
  \item Korngrößenverteilung beziehungsweise Sieblinie;
  \item Größtkorn;
  \item Feinanteile kleiner oder gleich $0{,}063\,\mathrm{mm}$;
  \item Kornform;
  \item Kornoberfläche;
  \item Anteil gebrochener oder gerundeter Körner.
\end{itemize}



Eine Korngruppe wird als Kornfraktion $d/D$ mit Unter- beziehungsweise
Überkorn bezeichnet. Ein Korngemisch ist ein natürliches Gemisch
$0/D$ oder ein Gemisch aus mehreren Korngruppen.


Dabei gilt:

\begin{itemize}
  \item $d$ ist der untere, eine Körnung beschreibende Durchmesser;
  \item $D$ ist der obere, eine Körnung beschreibende Durchmesser.
\end{itemize}



\section{Korngrößenverteilung und Sieblinie}

Die Korngrößenverteilung wird mit einem Siebverfahren nach
ÖNORM EN 933-1 bestimmt. Das Ergebnis wird als Sieblinie dargestellt.


Die Sieblinie zeigt den prozentuellen Durchgang einer Gesteinskörnung
durch unterschiedliche Siebe. Sie dient dazu, die Verteilung der
Korngrößen zu beurteilen und mit einem vorgegebenen
Sieblinienbereich zu vergleichen.


\begin{darstellungsbox}
\textbf{Darstellung: Sieblinie einer Gesteinskörnung}

Die Darstellung zeigt auf der horizontalen Achse die Siebgröße in
Millimetern. Auf der vertikalen Achse ist der Siebdurchgang in Prozent
dargestellt.

Eine rote Linie stellt die Sieblinie einer Gesteinskörnung dar. Blaue
Begrenzungslinien zeigen einen vorgegebenen zulässigen
Sieblinienbereich. Als Beispiel ist ein Korngemisch $0/32$ dargestellt.

Die dargestellten Siebgrößen reichen von $0{,}063\,\mathrm{mm}$ über
$0{,}2\,\mathrm{mm}$, $0{,}63\,\mathrm{mm}$, $2\,\mathrm{mm}$,
$6{,}3\,\mathrm{mm}$ und $20\,\mathrm{mm}$ bis
$63\,\mathrm{mm}$.
\end{darstellungsbox}

\section{Feinanteile}

Als Feinanteile werden Körner mit einer Korngröße von höchstens

\begin{equation}
  0{,}063\,\mathrm{mm}
\end{equation}

bezeichnet. 

Die Kategorie \texttt{f3} bedeutet, dass der Feinanteil kleiner oder
gleich $0{,}063\,\mathrm{mm}$ höchstens

\begin{equation}
  3\,\mathrm{M.\%}
\end{equation}

beträgt. 

\section{Kornform}

Die Kornform wird mit der Kornformkennzahl \texttt{SI} beschrieben.
Die Bestimmung erfolgt mit dem Shape Index nach ÖNORM EN 933-4.
Dabei wird der Anteil nicht-kubischer Körner mit einem Verhältnis von
Länge zu Breite von mehr als drei erfasst. 

Die Kategorie \texttt{SI20} bedeutet, dass der Masseanteil schlecht
geformter Körner höchstens

\begin{equation}
  20\,\%
\end{equation}

beträgt. 

\begin{merksatzbox}
\textbf{Merksatz:}
Die Kornform beeinflusst die Eigenschaften einer Gesteinskörnung. Mit
dem Shape Index wird insbesondere der Anteil nicht-kubischer Körner
bestimmt. Die Kategorie \texttt{SI20} begrenzt den Anteil schlecht
geformter Körner auf maximal $20\,\%$.

\end{merksatzbox}

\section{Bruchflächigkeit}

Die Bruchflächigkeit beschreibt den Anteil gebrochener und gerundeter
Körner einer Gesteinskörnung. Gebrochene Körner werden mit
\texttt{c} für \textit{crushed}, vollständig gebrochene Körner mit
\texttt{tc} für \textit{totally crushed} und vollständig gerundete
Körner mit \texttt{tr} für \textit{totally rounded} bezeichnet.


\begin{longtable}{
  >{\centering\arraybackslash}p{0.20\textwidth}
  >{\raggedright\arraybackslash}p{0.58\textwidth}
}
\toprule
\textbf{Kategorie} &
\textbf{Bedeutung} \\
\midrule
\endhead

\texttt{C100/0} &
Gesteinskörnung aus $100\,\%$ vollständig gebrochenen Körnern
\texttt{tc}. \\

\texttt{C90/1} &
Mindestens $90\,\%$ gebrochene Körner \texttt{c}; davon mindestens
$30\,\%$ vollständig gebrochene Körner \texttt{tc}; maximal
$1\,\%$ vollständig gerundete Körner \texttt{tr}. \\

\texttt{C50/30} &
Mindestens $50\,\%$ gebrochene Körner \texttt{c}; maximal
$30\,\%$ vollständig gerundete Körner \texttt{tr}. \\

\texttt{C50/10} &
Mindestens $50\,\%$ gebrochene Körner \texttt{c}; maximal
$10\,\%$ vollständig gerundete Körner \texttt{tr}. \\

\bottomrule
\end{longtable}



\section{Frost-Tau-Widerstand}

Der Frost-Tau-Widerstand von Gesteinskörnung wird nach
ÖNORM EN 1367-1 bestimmt. Dabei wird der Masseverlust einer
Gesteinskörnung nach zehn Frost-Tau-Wechseln ermittelt. Die
Durchführung erfolgt in Frostdosen. 

Der Frost-Tau-Wert $F$ wird mit folgender Gleichung berechnet:

\begin{equation}
  F =
  \frac{M_1-M_2}{M_1}
  \times 100
\end{equation}

Dabei bedeuten:

\begin{itemize}
  \item $M_1$: gesamte Ausgangstrockenmasse der drei
  Einzelmessproben in Gramm;
  \item $M_2$: gesamte Endtrockenmasse der drei Einzelmessproben,
  die auf dem festgelegten Analysensieb zurückgeblieben sind, in Gramm;
  \item $F$: prozentualer Masseverlust nach der
  Frost-Tau-Wechselbeanspruchung.
\end{itemize}



\begin{longtable}{
  >{\centering\arraybackslash}p{0.18\textwidth}
  >{\centering\arraybackslash}p{0.23\textwidth}
  >{\raggedright\arraybackslash}p{0.33\textwidth}
}
\toprule
\textbf{Kategorie} &
\textbf{Masseverlust} &
\textbf{Bedeutung} \\
\midrule
\endhead

F1 &
$\leq 1\,\%$ &
Frost-Tau-Widerstand \\

F2 &
$\leq 2\,\%$ &
Frost-Tau-Widerstand \\

F4 &
$\leq 4\,\%$ &
Frost-Tau-Widerstand \\

angegeben &
$>4\,\%$ &
Der Wert ist anzugeben. \\

FNR &
-- &
Keine Anforderung. \\

\bottomrule
\end{longtable}



\subsection{Frostbeständigkeit und Frostsicherheit}

Das Dokument unterscheidet zwischen Frostbeständigkeit und
Frostsicherheit:

\begin{itemize}
  \item Die \textbf{Frostbeständigkeit} betrifft das einzelne
  Gesteinskorn.

  \item Die \textbf{Frostsicherheit} betrifft das Gesteinsmaterial
  im eingebauten Zustand und hängt auch von der Korngrößenverteilung ab.
\end{itemize}



Für ungebundene Tragschichten werden Anforderungen insbesondere an
Tragfähigkeit und Frostsicherheit gestellt. Die Frostsicherheit wird
über Korn- und Mineralkriterien sowie durch Frosthebungsversuche
beurteilt. 

\begin{darstellungsbox}
\textbf{Darstellung: Frost-Tau-Prüfung}

Bei der Frost-Tau-Prüfung werden drei Einzelmessproben in Frostdosen
wiederholt Frost- und Tauwechseln ausgesetzt. Vor dem Versuch wird die
Ausgangstrockenmasse $M_1$ bestimmt.

Nach zehn Frost-Tau-Wechseln wird die Endtrockenmasse $M_2$ jener
Körnung bestimmt, die auf dem festgelegten Analysensieb
zurückgeblieben ist. Aus $M_1$ und $M_2$ wird der prozentuale
Masseverlust $F$ berechnet.
\end{darstellungsbox}

\section{Chemische Anforderungen}

Gesteinskörnungen dürfen keine Bestandteile enthalten, die die
Dauerhaftigkeit oder das Erstarrungs- und Erhärtungsverhalten von Beton
beeinträchtigen. 

Treiberscheinungen können durch Schwefel verursacht werden. Lösliche
Alkalisulfate wie Gips und Anhydrit sowie Sulfide werden als schädlich
bezeichnet. Nicht-lösliche Sulfate wie Baryt beziehungsweise
Schwerspat werden als unbedenklich angeführt. 

Chloride werden als stahlangreifende Stoffe genannt. Zusätzlich dürfen
keine Stoffe enthalten sein, welche das Erstarren oder Erhärten des
Betons verändern. Als Beispiel nennt das Dokument Zucker oder
zuckerähnliche Stoffe. 

\begin{longtable}{
  >{\raggedright\arraybackslash}p{0.32\textwidth}
  >{\raggedright\arraybackslash}p{0.47\textwidth}
}
\toprule
\textbf{Stoff oder Stoffgruppe} &
\textbf{Bewertung laut Dokument} \\
\midrule
\endhead

Lösliche Alkalisulfate, Gips, Anhydrit und Sulfide &
Schädlich; können Treiberscheinungen verursachen. \\

Nicht-lösliche Sulfate, beispielsweise Baryt oder Schwerspat &
Unbedenklich. \\

Chloride &
Stahlangreifende Stoffe. \\

Zucker oder zuckerähnliche Stoffe &
Können das Erstarrungs- und Erhärtungsverhalten des Betons verändern. \\

\bottomrule
\end{longtable}



\section{Wasseraufnahme und Wasserlöslichkeit}

Die Kategorie \texttt{WA24\textsubscript{1}} beschreibt die
Wasseraufnahme nach 24-stündiger Wasserbadlagerung. Sie bedeutet, dass
die Wasseraufnahme höchstens

\begin{equation}
  1\,\%
\end{equation}

beträgt. 

Die Kategorie \texttt{WS} ist ein Maß für die Wasserlöslichkeit von
Fremdfüller. 

\section{Widerstand gegen Zertrümmerung und Verschleiß}

Der Widerstand von Gesteinskörnung gegen Zertrümmerung und Verschleiß
kann mit unterschiedlichen Kategorien beschrieben werden.


\begin{longtable}{
  >{\centering\arraybackslash}p{0.22\textwidth}
  >{\raggedright\arraybackslash}p{0.54\textwidth}
}
\toprule
\textbf{Kategorie} &
\textbf{Bedeutung} \\
\midrule
\endhead

\texttt{LA25} &
Los-Angeles-Wert höchstens $25$; Maß für den Widerstand gegen
Zertrümmerung der Körnung $8/11\,\mathrm{mm}$. \\

\texttt{MDE15} &
Micro-Deval-Wert höchstens $15$; Maß für den Widerstand gegen
Verschleiß. \\

\texttt{SZ} &
Maß für den Widerstand gegen Zertrümmerung gemäß Schlagversuch. \\

\texttt{AAV} &
Maß für den Widerstand gegen Oberflächenabrieb. \\

\texttt{AN} &
Maß für den Widerstand gegen Abrieb durch Spikereifen. \\

\bottomrule
\end{longtable}



\section{Weitere Kategorien}

Das Dokument nennt weitere Kategorien zur Beschreibung von
Gesteinskörnungen und Füllern.

\begin{longtable}{
  >{\centering\arraybackslash}p{0.22\textwidth}
  >{\raggedright\arraybackslash}p{0.54\textwidth}
}
\toprule
\textbf{Kategorie} &
\textbf{Bedeutung} \\
\midrule
\endhead

\texttt{AS0,8} &
Anteil säurelöslicher Sulfate höchstens $0{,}8\,\%$. \\

\texttt{BN} &
Bitumenzahl; beschreibt die scheinbare Viskosität von Fremdfüller. \\

\texttt{CC80} &
Calcium-Carbonatgehalt eines Füllers mindestens $80\,\%$. \\

\texttt{PWS} &
Maß für den Polierwiderstand feiner Gesteinskörnungen nach
Wehner/Schulze. \\

\texttt{SBLA} &
Maß für den Widerstand gegen Sonnenbrand von Basalt. \\

\texttt{SC1} &
Muschelschalengehalt höchstens $10\,\%$. \\

\texttt{V3,5} &
Volumenzunahme von Stahlwerkschlacke höchstens $3{,}5\,\%$. \\

\texttt{V28/45} &
Maß für den Hohlraumgehalt von Fremdfüller. \\

\texttt{GA90} &
Korngemisch $0/D$ mit maximal $10\,\%$ Überkorn. \\

\texttt{Gc90/15} &
Grobe Gesteinskörnung mit maximal $10\,\%$ Überkorn und maximal
$15\,\%$ Unterkorn. \\

\texttt{GF85} &
Feine Gesteinskörnung mit $d=0$ und maximal $15\,\%$ Überkorn. \\

\texttt{GTC} &
Toleranz bei der Angabe der Korngrößenverteilung. \\

\texttt{Ka20} &
Calciumhydroxid-Gehalt mindestens $20\,\%$. \\

\texttt{MBF10} &
Methylenblaugehalt höchstens $10\,\mathrm{M.\%}$ zur Beurteilung der
Feinanteile kleiner oder gleich $0{,}063\,\mathrm{mm}$. \\

\bottomrule
\end{longtable}



\section{Zusammenfassung}

Gesteinskörnung muss geometrische, mechanische, physikalische und
chemische Anforderungen erfüllen. Wesentliche Merkmale sind
Korngrößenverteilung, Kornform, Bruchflächigkeit, Frost-Tau-Widerstand,
Wasseraufnahme, Widerstand gegen Zertrümmerung und Verschleiß sowie die
chemische Unbedenklichkeit.


\begin{merksatzbox}
\textbf{Merksatz:}
Für eine dauerhafte Verwendung in Beton, Asphalt oder Tragschichten
muss Gesteinskörnung unter anderem eine geeignete Korngrößenverteilung,
ausreichende Frost-Tau-Beständigkeit, geeignete Kornform und
chemische Unbedenklichkeit aufweisen.

\end{merksatzbox}
% =================================================
% KAPITEL 5: METALLISCHE WERKSTOFFE
% Diesen gesamten Block vor \end{document} einfügen.
% =================================================

\chapter{Metallische Werkstoffe}

\section{Kristalliner Aufbau von Metallen}

Die Eigenschaften von Metallen werden durch ihren besonderen atomaren
Aufbau bestimmt. Im festen Zustand ordnen sich die Atome in
regelmäßigen Raumgittern an. In diesen Gittern können sich die
Valenzelektronen nahezu frei bewegen. Diese Elektronenkonfiguration ist
entscheidend für die Leitfähigkeit und die Verformbarkeit von Metallen.


Abhängig vom Material, von den Randbedingungen und vom Energiezustand
können unterschiedliche Kristallstrukturen entstehen. Diese
Kristallstrukturen haben einen großen Einfluss auf die Eigenschaften
des Werkstoffs. 

\begin{merksatzbox}
\textbf{Merksatz:}
Die Eigenschaften metallischer Werkstoffe werden wesentlich durch ihren
atomaren Aufbau und die daraus entstehende Kristallstruktur bestimmt.
Die nahezu frei beweglichen Valenzelektronen sind insbesondere für
Leitfähigkeit und Verformbarkeit verantwortlich.

\end{merksatzbox}

\section{Wichtige Kristallstrukturen}

Das Dokument nennt drei wichtige Kristallstrukturen von Metallen:

\begin{itemize}
  \item das kubisch raumzentrierte Gitter, kurz \textbf{krz};
  \item das kubisch flächenzentrierte Gitter, kurz \textbf{kfz};
  \item das Gitter mit hexagonal dichtester Packung, kurz
  \textbf{hdp}.
\end{itemize}



\begin{longtable}{
  >{\raggedright\arraybackslash}p{0.34\textwidth}
  >{\raggedright\arraybackslash}p{0.42\textwidth}
}
\toprule
\textbf{Kristallstruktur} &
\textbf{Abkürzung} \\
\midrule
\endhead

Kubisch raumzentriertes Gitter &
krz \\

Kubisch flächenzentriertes Gitter &
kfz \\

Gitter mit hexagonal dichtester Packung &
hdp \\

\bottomrule
\end{longtable}



\section{Kristallstrukturen von Eisen}

Für Eisen werden im Dokument zwei Kristallstrukturen ausdrücklich
genannt:

\begin{itemize}
  \item $\alpha$-Eisen besitzt ein kubisch raumzentriertes Gitter;
  \item $\gamma$-Eisen besitzt ein kubisch flächenzentriertes Gitter.
\end{itemize}



\begin{longtable}{
  >{\raggedright\arraybackslash}p{0.30\textwidth}
  >{\raggedright\arraybackslash}p{0.45\textwidth}
}
\toprule
\textbf{Zustand des Eisens} &
\textbf{Kristallstruktur} \\
\midrule
\endhead

$\alpha$-Eisen &
kubisch raumzentriertes Gitter, krz \\

$\gamma$-Eisen &
kubisch flächenzentriertes Gitter, kfz \\

\bottomrule
\end{longtable}



\begin{darstellungsbox}
\textbf{Darstellung: Kristallgitter von Eisen}

Die Darstellung ordnet dem $\alpha$-Eisen das kubisch raumzentrierte
Gitter zu. Dabei befinden sich Atome an den Ecken eines Würfels sowie
ein weiteres Atom im Zentrum des Würfels.

Dem $\gamma$-Eisen wird das kubisch flächenzentrierte Gitter
zugeordnet. Dabei befinden sich Atome an den Ecken eines Würfels sowie
zusätzliche Atome in den Flächenmitten.

Die beiden Gitterformen verdeutlichen, dass sich die räumliche
Anordnung der Atome in einem metallischen Werkstoff ändern kann.
\end{darstellungsbox}

\section{Formgedächtnislegierungen}

Das Dokument behandelt Formgedächtnislegierungen im Zusammenhang mit
eisenbasierten Formgedächtnislegierungen. Diese können durch thermische
Nachbehandlung aktiviert werden. 

Das Materialverhalten einer Formgedächtnislegierung wird unter anderem
durch eine Abnahme der Materialsteifigkeit infolge einer
Metallgitterumwandlung beschrieben. Dabei wird eine Umwandlung von
Austenit zu Martensit genannt. 

\begin{merksatzbox}
\textbf{Merksatz:}
Bei Formgedächtnislegierungen kann sich die Materialsteifigkeit infolge
einer Metallgitterumwandlung verändern. Im Dokument wird dabei die
Umwandlung von Austenit zu Martensit genannt.

\end{merksatzbox}

\section{Einweg-Formgedächtnislegierung}

Das Dokument zeigt schematisch das Verhalten einer
Einweg-Formgedächtnislegierung. Dabei werden drei wesentliche Vorgänge
unterschieden:

\begin{enumerate}
  \item mechanisches Vordehnen der Legierung;
  \item Erwärmung der vorgedehnten Legierung bei freier
  Verformbarkeit;
  \item Erwärmung der vorgedehnten Legierung bei eingeschränkter
  beziehungsweise verhinderter freier Verformbarkeit.
\end{enumerate}



Bei steigender Belastung wird nach der Umwandlung des Gitters eine
elastische Verformung beschrieben. Bei einer Entlastung erfolgt eine
lineare Rückverformung. Gleichzeitig kann eine bleibende Dehnung aus
der Gitterumwandlung vorhanden sein. 

\begin{darstellungsbox}
\textbf{Darstellung: Verhalten einer Einweg-Formgedächtnislegierung}

Die Darstellung enthält Spannungs-Dehnungs- und
Spannungs-Temperatur-Verläufe einer Einweg-Formgedächtnislegierung.

Zunächst erfolgt eine mechanische Vordehnung. Danach wird die
vorgedehnte Legierung erwärmt. Die Darstellung unterscheidet dabei
einen Fall mit freier Verformbarkeit und einen Fall, bei dem die freie
Verformbarkeit eingeschränkt ist.

Als charakteristische Temperaturen sind $A_s$ und $A_f$ bezeichnet.
Die Darstellung zeigt außerdem die Rückverformung bei Entlastung und
die mit der Temperatur verknüpfte Spannungsentwicklung.
\end{darstellungsbox}

\section{Rückstellspannung}

Durch Abkühlen einer Formgedächtnislegierung kann die Spannung weiter
ansteigen. Die dabei erreichte Spannung wird als
\textbf{Rückstellspannung} bezeichnet. 

\begin{merksatzbox}
\textbf{Merksatz:}
Steigt die Spannung beim Abkühlen einer Formgedächtnislegierung weiter
an, wird die erreichte Spannung als Rückstellspannung bezeichnet.

\end{merksatzbox}

\section{Anwendung im Betonbau}

Das Dokument zeigt die Aktivierung einer Vorspannung durch
Formgedächtnislegierungen bei sofortigem Verbund. Dabei wird die
thermische Aktivierung der Formgedächtnislegierung im Zusammenhang mit
Vorspannung dargestellt. 

Die behandelte Quelle bezieht sich auf selbstvorspannende
ultrahochfeste Betonelemente, die mit eisenbasierten
Formgedächtnislegierungen bewehrt sind und durch thermische
Nachbehandlung aktiviert werden. 

\begin{darstellungsbox}
\textbf{Darstellung: Aktivierung einer Vorspannung}

Die Darstellung zeigt eine Formgedächtnislegierung im sofortigen
Verbund mit Beton. Durch thermische Aktivierung der
Formgedächtnislegierung kann eine Vorspannung aktiviert werden.

Die dargestellte Anwendung verbindet somit das temperaturabhängige
Materialverhalten der Legierung mit einer Anwendung in einem
Betonelement.
\end{darstellungsbox}

\section{Zusammenfassung}

Metalle besitzen im festen Zustand regelmäßig angeordnete Raumgitter.
Für Eisen nennt das Dokument $\alpha$-Eisen mit kubisch
raumzentriertem Gitter und $\gamma$-Eisen mit kubisch
flächenzentriertem Gitter. 

Formgedächtnislegierungen weisen ein spezielles, durch
Gitterumwandlungen geprägtes Materialverhalten auf. Sie können nach
mechanischem Vordehnen thermisch aktiviert werden. Beim Abkühlen kann
eine Rückstellspannung entstehen. \quelle{[2][^3]}

\begin{merksatzbox}
\textbf{Merksatz:}
Die im Dokument behandelten eisenbasierten Formgedächtnislegierungen
können durch thermische Nachbehandlung aktiviert werden. Beim Abkühlen
kann eine Rückstellspannung entstehen; dadurch ist unter anderem eine
Anwendung zur Aktivierung von Vorspannung möglich.
\quelle{[2][3]}
\end{merksatzbox}
% =================================================
% KAPITEL 6: GLAS
% =================================================
% =================================================
% KAPITEL 6: GLAS
% Diesen gesamten Block vor \end{document} einfügen.
% =================================================

\chapter{Glas}

\section{Struktur und Definition von Glas}

Glas ist ein anorganischer, meist silikatischer Festkörper. Es wird aus
Rohstoffen in einem Schmelzprozess hergestellt und erstarrt ohne
Kristallisation. 

Eine weitere Definition beschreibt Glas als anorganisches
Schmelzprodukt, das im Wesentlichen ohne Kristallisation erstarrt.


Siliziumdioxid kann mit kristallinem Aufbau als Quarz oder mit amorphem
Aufbau als Glas vorkommen. Glas unterscheidet sich damit von
kristallinen Stoffen durch seinen amorphen Aufbau. 

\begin{merksatzbox}
\textbf{Merksatz:}
Glas ist ein anorganischer, meist silikatischer Festkörper. Es entsteht
durch Schmelzen von Rohstoffen und erstarrt ohne Kristallisation.

\end{merksatzbox}

\section{Glaszustand}

Glas wird im physikalisch-chemischen Sinn als unterkühlte Flüssigkeit
bezeichnet. Gleichzeitig sind Gläser thermodynamisch metastabile
Festkörper. 

Unterhalb der Glasübergangstemperatur $T_g$ ist Glas spröde. Oberhalb
der Glasübergangstemperatur wird Glas plastisch verformbar.


Die Glasübergangstemperatur ist jene Temperatur, bei der die Viskosität
des Glases folgenden Wert erreicht:

\begin{equation}
  \eta = 10^{12}\,\mathrm{Pa \cdot s}
\end{equation}

Für übliche silikatische Gläser liegt der Transformationsbereich bei:

\begin{equation}
  520 \text{ bis } 550\,^\circ\mathrm{C}
\end{equation}



\begin{darstellungsbox}
\textbf{Darstellung: Übergang von kristallinen und glasigen Stoffen}

Kristalline Stoffe besitzen beim Übergang vom festen in den flüssigen
Zustand einen klar definierten Schmelzpunkt. Glasige Stoffe zeigen
dagegen keinen einzelnen scharfen Schmelzpunkt.

Bei Glas erfolgt der Übergang vom spröden Zustand zum plastisch
verformbaren beziehungsweise viskosen Zustand über einen breiteren
Temperaturbereich. Die Glasübergangstemperatur wird mit $T_g$
bezeichnet.
\end{darstellungsbox}

\section{Physikalische Eigenschaften}

Die mittlere Dichte von Natron-Kalk-Glas und Kali-Kalk-Glas beträgt:

\begin{equation}
  \rho \approx 2{,}5 \cdot 10^3\,\mathrm{kg/m^3}
\end{equation}

Bleiglas besitzt eine höhere Dichte zwischen:

\begin{equation}
  \rho =
  3{,}5 \text{ bis } 4{,}8 \cdot 10^3\,\mathrm{kg/m^3}
\end{equation}



Der Elastizitätsmodul von Glas hängt von der Zusammensetzung ab. Er
liegt im Bereich von:

\begin{equation}
  E =
  35\,000 \text{ bis } 75\,000\,\mathrm{N/mm^2}
\end{equation}

Der E-Modul ist außerdem stark temperaturabhängig. Bis zur
Transformationstemperatur nimmt er langsam ab; bei weiterer
Temperatursteigerung fällt er rasch gegen null. 

Die Härte von Glas beträgt auf der Mohs'schen Härteskala ungefähr:

\begin{equation}
  H = 6
\end{equation}

Glas ist damit ungefähr so hart wie Feldspat, jedoch weicher als Quarz
mit einer Mohs-Härte von $7$. Feiner Staub kann dennoch Kratzer auf
Glas hinterlassen. 

\begin{longtable}{
  >{\raggedright\arraybackslash}p{0.42\textwidth}
  >{\raggedright\arraybackslash}p{0.30\textwidth}
}
\toprule
\textbf{Eigenschaft} &
\textbf{Wert beziehungsweise Beschreibung} \\
\midrule
\endhead

Dichte von Natron-Kalk-Glas und Kali-Kalk-Glas &
ca.\ $2{,}5 \cdot 10^3\,\mathrm{kg/m^3}$ \\

Dichte von Bleiglas &
$3{,}5$ bis $4{,}8 \cdot 10^3\,\mathrm{kg/m^3}$ \\

Elastizitätsmodul &
$35\,000$ bis $75\,000\,\mathrm{N/mm^2}$ \\

Mohs-Härte &
ca.\ $6$ \\

Transformationsbereich silikatischer Gläser &
ca.\ $520$ bis $550\,^\circ\mathrm{C}$ \\

\bottomrule
\end{longtable}



\section{Thermische und optische Eigenschaften}

Der lineare thermische Ausdehnungskoeffizient von Bauglas beträgt rund:

\begin{equation}
  \alpha =
  9 \cdot 10^{-6}\,\mathrm{K^{-1}}
\end{equation}

Bei großflächigen Glasscheiben in Rahmenkonstruktionen müssen mögliche
Verformungsdifferenzen zwischen Rahmen und Glas berücksichtigt werden,
da Bruchgefahr bestehen kann. 

Die Wärmeleitzahl von Glas beträgt ungefähr:

\begin{equation}
  \lambda \approx 0{,}8\,\mathrm{W/(m \cdot K)}
\end{equation}

Glas ist damit ein relativ schlechter Wärmeleiter. 

Die Lichtdurchlässigkeit ist eine wesentliche Eigenschaft von Glas.
Sie erreicht jedoch nicht $100\,\%$. Die Lichtdurchlässigkeit
verschiedener Glasarten ist in der folgenden Tabelle dargestellt.

\begin{longtable}{
  >{\raggedright\arraybackslash}p{0.42\textwidth}
  >{\centering\arraybackslash}p{0.27\textwidth}
}
\toprule
\textbf{Glasart} &
\textbf{Lichtdurchlässigkeit} \\
\midrule
\endhead

Fensterglas &
$90$ bis $92\,\%$ \\

Kristallspiegelglas &
$90$ bis $92\,\%$ \\

Gartenglas &
$88$ bis $90\,\%$ \\

Drahtspiegelglas &
ca.\ $88\,\%$ \\

Gussglas und Drahtglas &
$80$ bis $90\,\%$ \\

Verbund-Sicherheitsglas &
$85$ bis $90\,\%$ \\

Mehrscheibenglas mit Metallbeschichtung &
$40$ bis $60\,\%$ \\

\bottomrule
\end{longtable}



Die Lichtdurchlässigkeit wird unter anderem durch Reflexion,
Lichtbrechung und Streuung beeinflusst. Handelsübliches Glas besitzt
eine Brechzahl zwischen ungefähr $1{,}3$ und $2{,}1$. Bei senkrechtem
Lichteinfall werden ungefähr $8\,\%$ des Lichts reflektiert.


\section{Mechanische Eigenschaften}

Glas ist im erkalteten Zustand ein spröder Werkstoff. Es kann
Spannungsspitzen nicht durch plastische Verformung abbauen.


Wird die Bruchdehnung von ungefähr

\begin{equation}
  \varepsilon_\mathrm{B} \approx 0{,}1\,\%
\end{equation}

erreicht, bricht Glas ohne Einschnürung. Die Bruchflächen sind
muschelförmig. 

Glas verhält sich bis zum Bruch linear-elastisch. Im Gegensatz zu Stahl
und Beton tritt bei Glas vor dem Bruch keine ausgeprägte plastische
Verformung auf. 

Die praktische Festigkeit von Glas liegt im konstruktiven Bereich
zwischen:

\begin{equation}
  30 \text{ bis } 90\,\mathrm{N/mm^2}
\end{equation}

Die theoretische Festigkeit ist deutlich höher. Der Unterschied wird
vor allem durch feine Mikrorisse an der Glasoberfläche verursacht.


\begin{longtable}{
  >{\raggedright\arraybackslash}p{0.42\textwidth}
  >{\centering\arraybackslash}p{0.28\textwidth}
}
\toprule
\textbf{Glaserzeugnis} &
\textbf{Charakteristische Biegezugfestigkeit} \\
\midrule
\endhead

Floatglas &
$45\,\mathrm{N/mm^2}$ \\

TVG aus Floatglas &
$70\,\mathrm{N/mm^2}$ \\

ESG aus Floatglas &
$120\,\mathrm{N/mm^2}$ \\

Chemisch vorgespanntes Glas &
$150\,\mathrm{N/mm^2}$ \\

\bottomrule
\end{longtable}



\begin{merksatzbox}
\textbf{Merksatz:}
Glas ist spröde und verhält sich bis zum Bruch linear-elastisch. Es
kann Spannungsspitzen nicht durch plastische Verformung abbauen.

\end{merksatzbox}

\section{Chemische Beständigkeit}

Glasoberflächen werden durch eine calciumoxid- und
siliziumdioxidreiche Oberflächenschicht vor saurem chemischem Angriff
geschützt. Wasser und die meisten Säuren greifen Glas daher praktisch
nicht an. 

Eine wesentliche Ausnahme ist Flusssäure. Sie kann Glas in kurzer Zeit
zerstören. Die chemische Reaktion lautet:

\begin{equation}
  \mathrm{
  4HF + SiO_2
  \Rightarrow
  SiF_4 + 2H_2O
  }
\end{equation}

Flusssäure wird daher zum Ätzen von Glas verwendet. 

Laugen, beispielsweise Kalkwasser oder Mörtelwasser, greifen Glas bei
dauernder Einwirkung an. Dabei wird Siliziumdioxid aus dem Glas
herausgelöst. Das Glas kann dadurch trüb werden und Streifen bilden.


\section{Rohstoffe für die Glasherstellung}

Die fein gemahlenen und möglichst reinen Rohstoffe werden entsprechend
ihren Massenanteilen gemischt. Dieses trockene Gemisch wird als
\textbf{Gemenge} bezeichnet. 

Die wichtigsten Grundbestandteile von Glas sind:

\begin{longtable}{
  >{\raggedright\arraybackslash}p{0.31\textwidth}
  >{\centering\arraybackslash}p{0.23\textwidth}
  >{\centering\arraybackslash}p{0.20\textwidth}
}
\toprule
\textbf{Rohstoff} &
\textbf{Chemisches Symbol} &
\textbf{Anteil} \\
\midrule
\endhead

Quarz als Hauptbestandteil &
$\mathrm{SiO_2}$ &
ca.\ $70\,\%$ \\

Soda als Flussmittel &
$\mathrm{Na_2CO_3}$ &
$11$ bis $16\,\%$ \\

Kalksteinmehl als Stabilisator &
$\mathrm{CaCO_3}$ &
$5$ bis $12\,\%$ \\

Magnesiumoxid &
$\mathrm{MgO}$ &
$0$ bis $5\,\%$ \\

Aluminiumoxid &
$\mathrm{Al_2O_3}$ &
$0{,}2$ bis $2\,\%$ \\

\bottomrule
\end{longtable}



\subsection{Quarzsand}

Siliziumdioxid in Form von Quarzsand ist der wichtigste Grundrohstoff
der Glasherstellung. Der Quarzsand muss eine hohe Reinheit besitzen.
Bereits Verunreinigungen von etwa $0{,}1\,\%$ Eisenoxid
$\mathrm{Fe_2O_3}$ können ihn für anspruchsvollere Anwendungen
unbrauchbar machen. 

Quarzsand ist der Hauptbestandteil des Gemenges und besitzt einen
Anteil von mindestens $60\,\%$. 

\subsection{Flussmittel}

Flussmittel werden eingesetzt, um die Schmelztemperatur abzusenken.
Dadurch wird der Schmelzprozess effizienter und wirtschaftlicher.


Glasscherben wirken als Flussmittel. Sie beschleunigen das
Flüssigwerden des Sandgemenges und können aus innerbetrieblichem
Ausschuss oder aus Recycling stammen. 

Soda ist Natriumcarbonat:

\begin{equation}
  \mathrm{Na_2CO_3}
\end{equation}

Beim Abkühlen geht Natriumoxid in das Glas über. Kohlendioxid entweicht
aus der Glasschmelze. 

\subsection{Stabilisatoren}

Kalk dient als Stabilisator. Bei Temperaturen von ungefähr
$1000\,^\circ\mathrm{C}$ entweicht aus Calciumcarbonat Kohlensäure.
Zurück bleibt Calciumoxid, das die chemischen und mechanischen
Eigenschaften des Glases verbessert. 

Aluminiumoxid wird dem Glasgemenge als alkali-haltiger Feldspat
zugegeben. Es erhöht die Zähigkeit des Glases im tieferen
Temperaturbereich. 

\section{Herstellungsprozess}

Der Glasherstellungsprozess gliedert sich in die Schritte:

\begin{enumerate}
  \item Mischen der Rohstoffe zum Gemenge,
  \item Schmelzen,
  \item Läutern,
  \item Formen,
  \item Kühlen.
\end{enumerate}



Das Gemenge wird im Ofen bei ungefähr

\begin{equation}
  1500\,^\circ\mathrm{C}
\end{equation}

eingeschmolzen. Durch Konvektion in der Schmelzwanne wird die
Glasschmelze homogenisiert. 

Die Formgebung kann durch Blasen, Ziehen, Walzen oder Gießen erfolgen.


Das Abkühlen muss langsam und stufenweise erfolgen. Die richtige
Abkühlgeschwindigkeit ist für die Qualität der Glaserzeugung von großer
Bedeutung. 

\begin{darstellungsbox}
\textbf{Darstellung: Glasherstellungsprozess}

Die Rohstoffe Quarz, Kalk und Soda werden zunächst gemischt. Die
fertige Mischung wird Gemenge genannt.

Das Gemenge wird bei ungefähr $1500\,^\circ\mathrm{C}$ eingeschmolzen.
Beim Läutern werden Blasen aus der Glasschmelze ausgetrieben.
Anschließend wird das Glas geformt, beispielsweise durch Blasen,
Ziehen, Walzen oder Gießen.

Nach der Formgebung folgt ein langsames und stufenweises Abkühlen.
\end{darstellungsbox}

\section{Glasarten}

Je nach Zusammensetzung und Einsatzbereich werden unterschiedliche
Glasarten verwendet. Die Eigenschaften hängen von der Zusammensetzung,
der Herstellung und der thermischen Vorgeschichte ab. 

\subsection{Natron-Kalk-Glas}

Natron-Kalk-Glas wird durch Zusammenschmelzen von Quarz, Soda und Kalk
hergestellt. Es wird insbesondere als Fensterglas, Tafelglas,
Flaschenglas und Spiegelglas verwendet. 

Die Reaktionsgleichung für die Glasbildung lautet:

\begin{equation}
  \mathrm{
  Na_2CO_3 + CaCO_3 + 6SiO_2
  \Rightarrow
  Na_2O \cdot CaO \cdot 6SiO_2 + 2CO_2
  }
\end{equation}



\subsection{Kali-Kalk-Glas}

Kali-Kalk-Glas wird mit Pottasche anstelle von Soda hergestellt. Es ist
auch als böhmisches Glas bekannt. Es ist schwerer schmelzbar, härter und
stärker lichtbrechend als Natronglas. 

Kali-Kalk-Glas wird für feinere und geschliffene Gegenstände verwendet.


\subsection{Bleiglas}

Bleiglas entsteht durch die Zugabe von Bleioxid. Es besitzt eine hohe
Dichte und ein starkes Lichtbrechungsvermögen. Es wird als
Bleikristallglas vor allem für geschliffene Glaswaren verwendet.


\subsection{Borosilikatglas}

Borosilikatglas entsteht, wenn ein Teil des Siliziumdioxids durch
Bortrioxid ersetzt wird. Der Anteil an Bortrioxid beträgt etwa
$8$ bis $13\,\%$ der Masse. 

Borosilikatglas besitzt einen geringeren Ausdehnungskoeffizienten und
ist daher weniger empfindlich gegenüber raschem Erhitzen und Abkühlen.
Es ist außerdem widerstandsfähiger gegen Wasser und Säuren.


Als Beispiele werden Jenaer Glas, Pyrexglas, Silexglas, Resistglas und
Duraxglas genannt. 

\subsection{Schaumglas}

Schaumglas entsteht durch gemeinsames Erhitzen von Glaspulver und einem
gasbildenden Stoff, beispielsweise Kohlenstoff. Es entsteht ein
geschlossenzelliger, spröd-harter und relativ schwerer Schaumstoff.


Schaumglas wird hauptsächlich in Plattenform eingesetzt, wenn neben
geringer Wärmeleitfähigkeit auch Unbrennbarkeit und Fäulnissicherheit
erforderlich sind. 

Der Anwendungstemperaturbereich liegt zwischen:

\begin{equation}
  -250\,^\circ\mathrm{C}
  \text{ und }
  +450\,^\circ\mathrm{C}
\end{equation}



\section{Floatglas}

Beim Floatglas-Verfahren wird die Glasschmelze auf ein flüssiges
Zinnbad geleitet. Da Glas eine geringere Dichte besitzt als flüssiges
Zinn, schwimmt die Glasschmelze auf dem Zinnbad. 

Durch die ebene Oberfläche des Zinnbads und die Oberflächenspannung
entsteht ein planparalleles Glasband. 

Die Glasschmelze besitzt beim Eintritt in das Floatbad ungefähr eine
Temperatur von:

\begin{equation}
  1100\,^\circ\mathrm{C}
\end{equation}

Danach wird das Glasband im Kühltunnel und auf Rollen planmäßig und
spannungsfrei auf Raumtemperatur abgekühlt. 

Die übliche Scheibendicke von Floatglas liegt zwischen:

\begin{equation}
  1{,}5 \text{ und } 19{,}0\,\mathrm{mm}
\end{equation}

In Ausnahmefällen können Scheiben bis zu $24{,}0\,\mathrm{mm}$ Dicke
hergestellt werden. 

\begin{darstellungsbox}
\textbf{Darstellung: Floatglas-Verfahren}

Die Glasschmelze wird auf ein flüssiges Zinnbad geleitet. Aufgrund der
geringeren Dichte schwimmt das Glas auf dem geschmolzenen Zinn.

Durch die plane Oberfläche des Zinnbads und die Oberflächenspannung
bildet sich ein ebenes, planparalleles Glasband. Dieses wird in einer
Kühlstraße beziehungsweise in einem Kühltunnel spannungsfrei
abgekühlt und anschließend zugeschnitten.
\end{darstellungsbox}

\section{Einscheibensicherheitsglas}

Einscheibensicherheitsglas wird mit \texttt{ESG} abgekürzt. Es handelt
sich um voll vorgespanntes Glas. Die Vorspannung wird durch einen
thermischen oder chemischen Prozess eingebracht. 

Beim thermischen Vorspannen werden Floatglasscheiben auf folgende
Temperatur erhitzt:

\begin{equation}
  620 \text{ bis } 650\,^\circ\mathrm{C}
\end{equation}

Danach werden sie durch Konvektion von der Oberfläche her rasch
abgekühlt. 

Durch die schnelle Abkühlung entsteht an der Glasoberfläche eine
Druckspannung. Im Inneren entsteht aus Gleichgewichtsgründen eine
Zugspannung. 

Die Druckspannung an der Oberfläche beträgt ungefähr:

\begin{equation}
  10{,}0\,\mathrm{kN/cm^2}
\end{equation}

Die Zugspannung im Inneren beträgt ungefähr:

\begin{equation}
  5{,}0\,\mathrm{kN/cm^2}
\end{equation}

Die Oberflächendruckspannung erhöht die Biegefestigkeit und die
Schlagfestigkeit. 

\subsection{Bruchverhalten von ESG}

Beim Bruch zerspringt ESG aufgrund der gespeicherten Energie aus dem
Eigenspannungszustand in kleine würfelförmige Bruchstücke. Diese
Bruchstücke besitzen eine Größe von weniger als:

\begin{equation}
  1\,\mathrm{cm^2}
\end{equation}

Dadurch wird das Risiko größerer Schnittverletzungen reduziert.


\begin{merksatzbox}
\textbf{Merksatz:}
ESG wird thermisch vorgespannt. Beim Bruch zerfällt es in kleine,
stumpfkantige Bruchstücke. Dadurch wird das Risiko schwerer
Schnittverletzungen gesenkt.

\end{merksatzbox}

\subsection{Spontanbruch von ESG}

Bei der Herstellung von Einscheibensicherheitsglas können
Nickel-Sulfid-Einschlüsse entstehen. Durch eine verzögerte
Phasenumwandlung von $\alpha$-NiS zu $\beta$-NiS kann sich das
Volumen um ungefähr $4\,\%$ vergrößern. 

Der Temperaturausdehnungskoeffizient von Nickel-Sulfid ist größer als
jener von Glas. Durch Temperaturerhöhung können innere Spannungen
entstehen. Überschreiten diese Spannungen die vorhandene Glasfestigkeit,
kann es zu einem Spontanbruch kommen. \quelle{[20][^4]}

Als Gegenmaßnahme wird ein Heißlagerungstest durchgeführt. Dabei wird
ESG auf folgende Temperatur erwärmt:

\begin{equation}
  290\,^\circ\mathrm{C}
\end{equation}

Die Dauer des Tests beträgt:

\begin{equation}
  8\,\mathrm{h}
\end{equation}

Dadurch wird das Wachstum von Nickel-Sulfid beschleunigt. Ein möglicher
Spontanbruch tritt dann bereits im Heißlagerungsofen auf.


\section{Teilvorgespanntes Glas}

Teilvorgespanntes Glas wird mit \texttt{TVG} abgekürzt. Es wird auf eine
geringere Temperatur als ESG erhitzt und langsamer sowie kontrollierter
abgekühlt. 

Die Druckvorspannung von TVG beträgt ungefähr:

\begin{equation}
  5{,}0\,\mathrm{kN/cm^2}
\end{equation}

TVG kann derzeit bis zu einer Dicke von ungefähr
$10{,}0\,\mathrm{mm}$ hergestellt werden. Seine Festigkeit liegt
zwischen jener von Floatglas und ESG. 

\section{Verbundsicherheitsglas}

Verbundsicherheitsglas wird mit \texttt{VSG} abgekürzt. Es besteht aus
zwei oder mehr Glasscheiben, die durch eine oder mehrere elastische,
hochreißfeste Kunststofffolien aus Polyvinylbutyral verbunden sind.


Die Dicke einer PVB-Folie beträgt:

\begin{equation}
  0{,}38\,\mathrm{mm}
\end{equation}



Die Glasscheiben können aus Floatglas, ESG oder TVG bestehen.


Die Folien werden auf die Glasscheiben gewalzt. Anschließend wird der
Verbund im Autoklav hergestellt. Die Herstellung erfolgt bei ungefähr:

\begin{align}
  T &\approx 140\,^\circ\mathrm{C}, \\
  p &\approx 12\,\mathrm{bar},
\end{align}

über eine Dauer von:

\begin{equation}
  6 \text{ bis } 8\,\mathrm{h}
\end{equation}



Die Folie hält beim Bruch die Bruchstücke an der Folie fest. Bei
Verbundsicherheitsglas aus Floatglas oder TVG kann die Resttragfähigkeit
auch im gebrochenen Zustand relativ hoch sein. 

\section{Isolierglas}

Isolierglas besteht aus zwei oder mehreren Scheiben, die durch
Gasschichten getrennt sind. Als Füllgase werden Argon, Krypton und
Xenon genannt. Auch Luftschichten können verwendet werden.


Die Gas- oder Luftschicht dient der Wärmedämmung. Die Ränder der
gesamten Scheibe müssen gasdicht abgeschlossen sein. 

Der Abstand zwischen den Scheiben beträgt:

\begin{equation}
  7{,}0 \text{ bis } 20{,}0\,\mathrm{mm}
\end{equation}



Die einzelnen Scheiben können aus Floatglas, ESG, TVG oder VSG
bestehen. 

Bei der Bemessung muss die klimatisch bedingte Druckdifferenz des
eingeschlossenen Gases berücksichtigt werden. Diese entsteht durch
Änderungen von Temperatur, Luftdruck oder Ortshöhe. 

\begin{darstellungsbox}
\textbf{Darstellung: Isolierglas und Klimalast}

Isolierglas besteht aus mindestens zwei Glasscheiben. Zwischen den
Scheiben befindet sich eine Gas- oder Luftschicht. Abstandshalter,
Trockenmittel und Dichtungen bilden den Randverbund.

Ändern sich Temperatur, Luftdruck oder Ortshöhe, verändert sich der
Druck des eingeschlossenen Gases. Diese Einwirkung wird als
Klimalast bezeichnet und kann zu Verformungen der Glasscheiben führen.
\end{darstellungsbox}

\section{Bearbeitung und Beschichtung}

Die normale Schnittkante ist die einfachste Form der Kantenbearbeitung.
Sie wird dort verwendet, wo die Schnittkante im Rahmen liegt und keine
Verletzungsgefahr durch die scharfe Kante besteht. 

Weitere Bearbeitungsmöglichkeiten sind Schleifen und Polieren. Bei ESG
wird die Kante leicht gebrochen beziehungsweise abgeschrägt, damit die
scharfe Kante entschärft wird. 

Als Oberflächenbehandlungen werden genannt:

\begin{itemize}
  \item Emailierung,
  \item Ätzung,
  \item Sandstrahlung.
\end{itemize}



Beschichtungen können auf der Außenseite, der Innenseite oder im
Scheibenzwischenraum angebracht werden. Sie dienen insbesondere für:

\begin{itemize}
  \item Wärmeschutzglas,
  \item Sonnenschutzglas,
  \item Reflexionsschichten,
  \item vereinfachte Reinigung mit Lotus-Effekt.
\end{itemize}



\section{Zusammenfassung}

Glas ist ein amorpher, silikatischer Werkstoff, der ohne
Kristallisation erstarrt. Unterhalb der Glasübergangstemperatur ist es
spröde, oberhalb davon plastisch verformbar. 

Für Bauglas sind insbesondere Floatglas, ESG, TVG, VSG und Isolierglas
von Bedeutung. Floatglas bildet die Ausgangsbasis für zahlreiche
weiterverarbeitete Glasprodukte. 

\begin{merksatzbox}
\textbf{Merksatz:}
Floatglas entsteht auf einem flüssigen Zinnbad. ESG wird durch
thermisches Vorspannen verfestigt. VSG besteht aus mehreren
Glasscheiben mit PVB-Folien. Isolierglas besteht aus mehreren durch
Gas- oder Luftschichten getrennten Scheiben.
\quelle{[19][4]}
\end{merksatzbox}

% =================================================
% KAPITEL 7: HOLZ UND HOLZWERKSTOFFE
% Diesen Block vor \end{document} einfügen.
% =================================================

% =================================================
% KAPITEL 7: HOLZ UND HOLZWERKSTOFFE
% Diesen gesamten Block vor \end{document} einfügen.
% =================================================

\chapter{Holz und Holzwerkstoffe}

\section{Allgemeines}

Holz ist ein natürlicher, organischer Baustoff und zählt gemeinsam mit
Naturstein zu den ältesten Baumaterialien der Erde. Bereits in früher
Zeit wurde Holz für Pfahlbauten verwendet. Später kam es bei
Fachwerkbauten zum Einsatz. 

Heute wird Holz sowohl im klassischen Hochbau als auch für tragende
Konstruktionen verwendet. Zu den im Dokument genannten
Anwendungsgebieten zählen:

\begin{itemize}
  \item Dachstühle,
  \item Dachausbauten,
  \item Stiegen,
  \item Fertighäuser,
  \item Blockhäuser,
  \item Hallenbauten,
  \item Brückenkonstruktionen,
  \item Betonschalungen,
  \item Gerüstbau,
  \item Pilotierungen,
  \item kleinere Schutzbauten gegen Naturgefahren.
\end{itemize}



\subsection{Vorteile von Holz}

Holz ist ein nachwachsender heimischer Rohstoff. Es ist leicht zu
bearbeiten und besitzt gute bauphysikalische Eigenschaften.
Insbesondere in wärmetechnischer Hinsicht besitzt Holz günstige
Eigenschaften. 

Im Verhältnis zu seiner Festigkeit ist Holz ein leichter Baustoff.
Zusätzlich weist Holz eine hohe Biege- und Druckfestigkeit auf.


\begin{longtable}{
  >{\raggedright\arraybackslash}p{0.35\textwidth}
  >{\raggedright\arraybackslash}p{0.45\textwidth}
}
\toprule
\textbf{Vorteil} &
\textbf{Beschreibung} \\
\midrule
\endhead

Nachwachsender Rohstoff &
Holz wird als heimischer und nachwachsender Rohstoff beschrieben. \\

Geringes Gewicht &
Holz ist im Verhältnis zu seiner Festigkeit ein leichter Baustoff. \\

Bearbeitbarkeit &
Holz lässt sich leicht bearbeiten. \\

Mechanische Eigenschaften &
Holz weist eine hohe Biege- und Druckfestigkeit auf. \\

Bauphysikalische Eigenschaften &
Holz besitzt insbesondere günstige wärmetechnische Eigenschaften. \\

\bottomrule
\end{longtable}



\subsection{Nachteile von Holz}

Holz schwindet und quillt. Bei schwankender Luftfeuchtigkeit arbeitet
unzureichend getrocknetes Holz. Bei zu schneller Trocknung kann es zu
Rissbildungen kommen. 

Da Holz ein natürlich wachsender Rohstoff ist, sind seine Abmessungen
begrenzt und seine Festigkeitseigenschaften inhomogen. Holz ist
außerdem brennbar sowie anfällig gegenüber Fäulnis und Zerstörung durch
Insekten. 

\begin{longtable}{
  >{\raggedright\arraybackslash}p{0.35\textwidth}
  >{\raggedright\arraybackslash}p{0.45\textwidth}
}
\toprule
\textbf{Nachteil} &
\textbf{Beschreibung} \\
\midrule
\endhead

Quellen und Schwinden &
Holz verändert sein Volumen bei Änderung des Feuchtegehalts. \\

Rissbildung &
Bei zu schneller Trocknung können Risse entstehen. \\

Begrenzte Abmessungen &
Holz ist ein natürlich wachsender Rohstoff; seine Abmessungen sind
begrenzt. \\

Inhomogene Festigkeit &
Die Festigkeitseigenschaften können unterschiedlich sein. \\

Brennbarkeit &
Holz ist ein brennbarer Baustoff. \\

Biologische Anfälligkeit &
Ohne geeignete Maßnahmen kann Holz durch Fäulnis, Pilze und Insekten
geschädigt werden. \\

\bottomrule
\end{longtable}



\begin{merksatzbox}
\textbf{Merksatz:}
Holz besitzt ein günstiges Verhältnis von Festigkeit zu Gewicht und
gute wärmetechnische Eigenschaften. Nachteile sind insbesondere
Feuchtebedingte Volumenänderungen, Rissbildung, Brennbarkeit und die
Anfälligkeit gegenüber biologischen Schädigungen.

\end{merksatzbox}

\section{Holzfeuchte}

Die Holzfeuchte ist eines der wichtigsten Beschreibungsmerkmale von
Holz. Sie beeinflusst unter anderem Festigkeit, Verformbarkeit und
Masse. 

Der Feuchtegehalt beschreibt den prozentuellen Anteil des enthaltenen
Wassers bezogen auf die darrtrockene Masse:

\begin{equation}
  u =
  \frac{m_u - m_{\mathrm{dtr}}}
  {m_{\mathrm{dtr}}}
  \cdot 100
  \,[\%]
\end{equation}

Dabei bedeuten:

\begin{itemize}
  \item $u$: Feuchtegehalt in Prozent,
  \item $m_u$: Masse des Holzes im feuchten Zustand,
  \item $m_{\mathrm{dtr}}$: Masse des Holzes im darrtrockenen Zustand.
\end{itemize}



\subsection{Grenzzustände des Systems Holz--Wasser}

Das Dokument unterscheidet darrtrockenes Holz, Fasersättigung und
Wassersättigung.

\begin{longtable}{
  >{\raggedright\arraybackslash}p{0.28\textwidth}
  >{\raggedright\arraybackslash}p{0.52\textwidth}
}
\toprule
\textbf{Zustand} &
\textbf{Beschreibung} \\
\midrule
\endhead

Darrtrocken &
Es ist kein Wasser im Holz vorhanden. Die Holzfeuchte beträgt
$0\,\%$. \\

Fasersättigung &
Das Mikrosystem der Zellwand ist vollständig mit Wasser gefüllt. Die
Holzfeuchte liegt ungefähr bei $28$ bis $32\,\%$. \\

Wassersättigung &
Mikro- und Makrosystem des Holzes sind maximal mit Wasser gefüllt. Die
Holzfeuchte kann je nach Holzdichte zwischen $31\,\%$ bei Pockholz und
$770\,\%$ bei Balsa liegen. \\

\bottomrule
\end{longtable}



Zwischen darrtrockenem und fasergesättigtem Zustand stellt sich ein
Feuchtegleichgewicht ein. Dabei passt sich die Holzfeuchte an die
relative Luftfeuchtigkeit an. Dieser Bereich ist von der
Lufttemperatur abhängig. 

\subsection{Quellen und Schwinden}

Wenn Holz Feuchtigkeit verliert, verkleinert sich sein Volumen. Dieser
Vorgang wird als Schwinden bezeichnet. Nimmt Holz Wasser auf, vergrößert
sich sein Volumen; dies wird Quellen genannt. 

Das Schwinden beginnt, sobald der Feuchtegehalt unter die Fasersättigung
sinkt. Dies liegt bei einer Holzfeuchte von ungefähr weniger als
$30\,\%$. Starkes Schwinden kann zu Rissbildungen führen.


Holz weist in unterschiedlichen Schnittrichtungen ein unterschiedlich
starkes Quell- und Schwindverhalten auf. Das Quellen und Schwinden ist
in Tangentialrichtung größer als in Radialrichtung. Das Werfen des
Holzes wird auf dieses unterschiedliche Verhalten zurückgeführt.


\begin{darstellungsbox}
\textbf{Darstellung: Quellen und Schwinden von Holz}

Bei abnehmender Holzfeuchte verkleinert sich das Holzvolumen. Dieser
Vorgang wird Schwinden genannt. Bei zunehmender Holzfeuchte vergrößert
sich das Holzvolumen; dieser Vorgang wird Quellen genannt.

Das Quell- und Schwindverhalten ist richtungsabhängig. In
Tangentialrichtung treten stärkere Verformungen auf als in
Radialrichtung. Dadurch kann sich Holz verformen oder werfen.
\end{darstellungsbox}

\section{Mechanische Eigenschaften}

Die maßgebliche Kenngröße der mechanischen Eigenschaften ist die
Festigkeit. Sie ist als jene Grenzspannung definiert, bei der ein
Prüfkörper unter Belastung versagt. 

Das Dokument unterscheidet folgende Festigkeitsarten:

\begin{itemize}
  \item Zugfestigkeit,
  \item Druckfestigkeit,
  \item Biegefestigkeit,
  \item Scherfestigkeit,
  \item Spaltfestigkeit,
  \item Torsionsfestigkeit.
\end{itemize}



Für die Bemessung von Holzbauteilen sind Zug-, Druck- und
Biegefestigkeit besonders wichtig. Holz weist in Faserrichtung und quer
zur Faserrichtung stark unterschiedliche Festigkeiten auf.


\subsection{Zugfestigkeit}

Die Zugfestigkeit wird als Bruchkraft bezogen auf die
Querschnittsfläche definiert:

\begin{equation}
  \sigma_{\mathrm{ZB}} =
  \frac{F_{\mathrm{max}}}{A}
\end{equation}

Dabei bedeuten:

\begin{itemize}
  \item $\sigma_{\mathrm{ZB}}$: Zugfestigkeit in
  $\mathrm{N/mm^2}$,
  \item $F_{\mathrm{max}}$: Bruchkraft in $\mathrm{N}$,
  \item $A$: Querschnittsfläche in $\mathrm{mm^2}$.
\end{itemize}



Die Zugfestigkeit senkrecht zur Faserrichtung beträgt lediglich etwa
$5$ bis $10\,\%$ der Zugfestigkeit in Faserrichtung. Querzug-
beanspruchungen sollen im Holzbau daher möglichst vermieden werden.


\begin{merksatzbox}
\textbf{Merksatz:}
Die Zugfestigkeit von Holz senkrecht zur Faser beträgt nur ungefähr
$5$ bis $10\,\%$ der Zugfestigkeit in Faserrichtung. Querzug-
beanspruchungen sind im Holzbau daher weitgehend zu vermeiden.

\end{merksatzbox}

\section{Bestand und Nutzung in Österreich}

In Österreich sind ungefähr $48\,\%$ der Staatsfläche Wald. Dies
entspricht einer Waldfläche von ungefähr $40\,000\,\mathrm{km^2}$.


Der österreichische Waldbestand setzt sich laut Dokument zusammen aus:

\begin{itemize}
  \item $60{,}5\,\%$ Nadelholz,
  \item $25{,}3\,\%$ Laubholz,
  \item $14{,}2\,\%$ sonstigem Holz.
\end{itemize}



Fichte ist mit ungefähr $46{,}2\,\%$ der Gesamtwaldfläche die
dominierende Baumart. Die Rotbuche ist mit einem Anteil von
$10{,}5\,\%$ der am stärksten vertretene Laubbaum.


Der jährliche Holzzuwachs im österreichischen Ertragswald beträgt laut
Dokument ungefähr:

\begin{equation}
  29{,}2\,\text{Mio.\ Vorratsfestmeter}
\end{equation}



\section{Struktur und Aufbau von Holz}

Holz wird grundsätzlich in Laubholz und Nadelholz unterschieden. Zur
Bestimmung der Holzart und ihrer Eigenschaften werden anatomische
Unterschiede anhand unterschiedlicher Schnittrichtungen untersucht.


Die wesentlichen Schnittrichtungen sind:

\begin{itemize}
  \item Querschnitt,
  \item Tangentialschnitt,
  \item Radialschnitt.
\end{itemize}



\subsection{Anatomie des Stammes}

Im Zentrum des Stammes verläuft in Stammlängsrichtung der Markstrang.
Dieser ist in frühen Stadien für die Nährstoffversorgung verantwortlich.


Um das Mark herum liegen die Jahresringe. Ein Jahresring besteht aus
Frühholz und Spätholz und beschreibt den jährlichen Zuwachs des
Baumes. 

Frühholz entsteht im Frühjahr und dient vor allem der Wasserversorgung.
Spätholz bildet sich in den Sommer- und Herbstmonaten, erhöht die
Festigkeit und erscheint aufgrund von Dichteunterschieden dunkler.


Außen am Holz liegt das Kambium. Dort erfolgt durch Zellteilung das
Dickenwachstum des Baumes. Außen an das Kambium grenzt der Bast, auch
Innenrinde genannt. Die Borke schützt den Stamm vor mechanischen,
tierischen und menschlichen Einflüssen. 

Holzstrahlen verlaufen in radialer Richtung. Sie dienen dem Transport
und der Speicherung von Nährstoffen. 

\begin{longtable}{
  >{\raggedright\arraybackslash}p{0.27\textwidth}
  >{\raggedright\arraybackslash}p{0.53\textwidth}
}
\toprule
\textbf{Bestandteil} &
\textbf{Funktion oder Beschreibung} \\
\midrule
\endhead

Markstrang &
Verläuft im Zentrum des Stammes; ist in frühen Entwicklungsstadien für
die Nährstoffversorgung verantwortlich. \\

Frühholz &
Entsteht im Frühjahr und dient der Wasserversorgung. \\

Spätholz &
Entsteht im Sommer und Herbst; ist festigkeitssteigernd und erscheint
dunkler. \\

Kambium &
Bereich des Dickenwachstums durch Zellteilung. \\

Bast &
Innenrinde; liegt außen am Kambium. \\

Borke &
Schützt den Stamm und besteht aus abgestorbenem Bastgewebe. \\

Holzstrahlen &
Dienen dem radialen Transport und der Speicherung von Nährstoffen. \\

\bottomrule
\end{longtable}



\subsection{Chemische Zusammensetzung}

Holz ist ein natürlicher organischer Werkstoff mit einer komplexen
Struktur aus Makromolekülen. Der Wassergehalt kann zwischen
$15$ und $50\,\%$ schwanken. Ohne Wasser gliedert sich Holz laut
Dokument in:

\begin{itemize}
  \item Zellulose und Hemizellulose: $70\,\%$,
  \item Lignin: $25\,\%$,
  \item Extraktstoffe: $5\,\%$.
\end{itemize}



Zellulose, Hemizellulose und begleitende Kohlenhydrate bilden die
Gerüstsubstanz des Holzes. Lignin bildet die Kittsubstanz zwischen
diesen Bestandteilen. 

Extraktstoffe können beispielsweise Harze, Fette, Wachse, Farb- und
Gerbstoffe sein. Sie können chemische, biologische und physikalische
Eigenschaften wie Lichtbeständigkeit, Brandverhalten oder
Schädlingsresistenz beeinflussen. 

\begin{longtable}{
  >{\raggedright\arraybackslash}p{0.30\textwidth}
  >{\centering\arraybackslash}p{0.23\textwidth}
  >{\raggedright\arraybackslash}p{0.28\textwidth}
}
\toprule
\textbf{Bestandteil} &
\textbf{Anteil ohne Wasser} &
\textbf{Bedeutung} \\
\midrule
\endhead

Zellulose und Hemizellulose &
ca.\ $70\,\%$ &
Gerüstsubstanz des Holzes. \\

Lignin &
ca.\ $25\,\%$ &
Kittsubstanz zwischen den Gerüstbestandteilen. \\

Extraktstoffe &
ca.\ $5\,\%$ &
Beispielsweise Harze, Fette, Wachse, Farb- und Gerbstoffe. \\

\bottomrule
\end{longtable}



Die wichtigsten chemischen Elemente des Holzes sind Kohlenstoff,
Sauerstoff, Wasserstoff und Stickstoff.

\begin{longtable}{
  >{\raggedright\arraybackslash}p{0.30\textwidth}
  >{\centering\arraybackslash}p{0.18\textwidth}
  >{\centering\arraybackslash}p{0.22\textwidth}
}
\toprule
\textbf{Element} &
\textbf{Chemisches Zeichen} &
\textbf{Gewichtsanteil} \\
\midrule
\endhead

Kohlenstoff &
C &
ca.\ $50\,\%$ \\

Sauerstoff &
O &
ca.\ $43\,\%$ \\

Wasserstoff &
H &
ca.\ $6\,\%$ \\

Stickstoff &
N &
ca.\ $0{,}6\,\%$ \\

Mineralstoffe &
-- &
ca.\ $0{,}2\,\%$ \\

\bottomrule
\end{longtable}



\section{Massivholz und Bauholzarten}

Für Bauholz werden in Österreich sowohl Nadelhölzer als auch Laubhölzer
verwendet. Häufig eingesetzte Holzarten sind Fichte, Kiefer, Lärche,
Eiche und Buche. 

\subsection{Nadelhölzer}

Fichte ist gemessen an ihrem Anteil im Ertragswald die wichtigste
Baumart in Österreich. Fichtenholz ist leicht, weich und mäßig
schwindend. Seine Darrdichte beträgt ungefähr:

\begin{equation}
  \rho_{\mathrm{dtr}} \approx 410\,\mathrm{kg/m^3}
\end{equation}

Fichte weist eine relativ hohe Druck- und Biegefestigkeit auf, ist
jedoch gegenüber Fäulnis und Schädlingen anfällig. 

Kiefernholz besitzt eine Darrdichte von ungefähr:

\begin{equation}
  \rho_{\mathrm{dtr}} \approx 530\,\mathrm{kg/m^3}
\end{equation}

Es besitzt hinsichtlich der Widerstandsfähigkeit Vorteile gegenüber
Fichte. 

Lärchenholz besitzt eine Darrdichte von ungefähr:

\begin{equation}
  \rho_{\mathrm{dtr}} \approx 550\,\mathrm{kg/m^3}
\end{equation}

Lärche besitzt gute Festigkeitseigenschaften. Aufgrund ihres
Harzreichtums gilt sie als dauerhaft und gegenüber Schädlingen relativ
beständig. 

\subsection{Laubhölzer}

Eiche zählt zu den schweren und harten Hölzern. Ihre Darrdichte beträgt
ungefähr:

\begin{equation}
  \rho_{\mathrm{dtr}} \approx 670\,\mathrm{kg/m^3}
\end{equation}

Die Dauerhaftigkeit von Eichenholz wird auf den hohen Gerbsäuregehalt
und einen pH-Wert von ungefähr $3{,}5$ zurückgeführt. 

Buche ist ebenfalls ein schweres und hartes Holz. Ihre Darrdichte
beträgt ungefähr:

\begin{equation}
  \rho_{\mathrm{dtr}} \approx 680\,\mathrm{kg/m^3}
\end{equation}

Buchenholz schwindet stark und ist eher fäulnisempfindlich, besitzt
jedoch höhere Festigkeiten als Fichte. 

\begin{longtable}{
  >{\raggedright\arraybackslash}p{0.24\textwidth}
  >{\centering\arraybackslash}p{0.23\textwidth}
  >{\raggedright\arraybackslash}p{0.35\textwidth}
}
\toprule
\textbf{Holzart} &
\textbf{Darrdichte} &
\textbf{Ausgewählte Eigenschaften} \\
\midrule
\endhead

Fichte &
ca.\ $410\,\mathrm{kg/m^3}$ &
Leicht, weich, mäßig schwindend; relativ hohe Druck- und
Biegefestigkeit. \\

Kiefer &
ca.\ $530\,\mathrm{kg/m^3}$ &
Harzreich; gegenüber Fichte widerstandsfähiger. \\

Lärche &
ca.\ $550\,\mathrm{kg/m^3}$ &
Gute Festigkeit; harzreich, dauerhaft und relativ schädlingsbeständig. \\

Eiche &
ca.\ $670\,\mathrm{kg/m^3}$ &
Schwer, hart und dauerhaft. \\

Buche &
ca.\ $680\,\mathrm{kg/m^3}$ &
Schwer und hart; schwindet stark und ist eher fäulnisempfindlich. \\

\bottomrule
\end{longtable}



\section{Holzfehler}

Als natürlich wachsender Rohstoff ist Holz während des Wachstums
verschiedenen Beanspruchungen und Einwirkungen ausgesetzt. Dadurch
entstehen Holzfehler, die den Gebrauchswert vermindern oder Holz für
bestimmte Anwendungen unbrauchbar machen können. 

\subsection{Fehler der Stammform}

Baumart, Klima und Standort beeinflussen die Stammform. Ein möglichst
vollholziger und sich nur allmählich verjüngender Stamm wird als
günstig angesehen. 

Das Dokument nennt insbesondere:

\begin{itemize}
  \item Drehwuchs,
  \item Abholzigkeit,
  \item Krummschäftigkeit,
  \item Zwieselwuchs.
\end{itemize}



Drehwuchs kann bei stärkerer Ausprägung zur Nichteignung als Bauholz
führen. Krummschäftigkeit bewirkt eine bedeutende Wertminderung.


\subsection{Fehler des Stammquerschnitts}

Ein gleichmäßiger und runder Stammquerschnitt ist erwünscht. Häufige
Fehler des Stammquerschnitts sind:

\begin{itemize}
  \item exzentrischer Wuchs,
  \item Rotkernigkeit,
  \item spanrückiger Wuchs,
  \item Zwieselbildung.
\end{itemize}



\subsection{Ästigkeit}

Der Astanteil eines Baumes ist stark vom Standort abhängig. In dichten
Beständen können relativ hohe astfreie Stämme entstehen, da untere Äste
aufgrund von Lichtmangel absterben. 

Das Dokument unterscheidet weiße und schwarze Äste:

\begin{itemize}
  \item Weiße Äste sind zum Fällzeitpunkt noch lebendig und fest mit
  dem Stammholz verwachsen.

  \item Schwarze Äste sind abgestorben und durch das Dickenwachstum in
  das Stammholz eingewachsen.
\end{itemize}



Werden abgestorbene eingewachsene Äste während der Verarbeitung
angeschnitten, können sie herausfallen. Diese Durchfalläste sind nicht
erwünscht. 

\subsection{Weitere Holzfehler}

Fehler infolge äußerer Einwirkung sind beispielsweise Wundschäden,
Frostrisse, Schwindrisse, Risse in den Jahresringen und Kernrisse.


Holzfehler können auch durch Schädlingsbefall entstehen. Insekten können
Fraßgänge verursachen; Pilzbefall kann Braunfäule oder Weißfäule
hervorrufen. 

\section{Sortierung und Festigkeitsklassen}

Die Sortierung von Vollholz erfolgt hinsichtlich der Tragfähigkeit und
deren Einflussfaktoren. Sie kann visuell oder maschinell erfolgen.


Visuell sortiertes Holz wird in Sortierklassen mit dem Kürzel
\texttt{S} eingeordnet. Maschinell sortiertes Holz wird mit dem Kürzel
\texttt{MS} bezeichnet. 

Beispiele für Sortiermerkmale sind:

\begin{itemize}
  \item Äste,
  \item Jahrringbreite,
  \item Risse,
  \item Faserneigung,
  \item Verfärbung,
  \item Insektenbefall,
  \item Rohdichte,
  \item Holzfeuchtigkeit.
\end{itemize}



Die Festigkeitsklassen für Bauholz werden nach EN~338 angegeben. Für
Nadelholz nennt das Dokument beispielsweise die Klassen \texttt{C16},
\texttt{C24}, \texttt{C30}, \texttt{C35} und \texttt{C40}.


Für Laubholz werden unter anderem die Klassen \texttt{D30},
\texttt{D35}, \texttt{D40}, \texttt{D50}, \texttt{D60} und
\texttt{D70} genannt. 

\begin{merksatzbox}
\textbf{Merksatz:}
Holz ist ein Naturprodukt mit unterschiedlichen Eigenschaften. Deshalb
wird Bauholz visuell oder maschinell sortiert und einer
Festigkeitsklasse nach EN~338 zugeordnet.

\end{merksatzbox}

\section{Holzwerkstoffe}

Holzwerkstoffe sind weiterverarbeitete Produkte, deren
Materialeigenschaften während der Produktion beeinflusst werden können.


Das Dokument unterscheidet folgende Gruppen:

\begin{itemize}
  \item Werkstoffe auf Vollholzbasis,
  \item Werkstoffe auf Furnierbasis,
  \item Werkstoffe auf Spanbasis,
  \item Werkstoffe auf Faserbasis,
  \item Verbundwerkstoffe.
\end{itemize}



\subsection{Werkstoffe auf Vollholzbasis}

Wichtige Vertreter der Werkstoffe auf Vollholzbasis sind
Brettschichtholz und Massivholzplatten. 

Brettschichtholz besteht aus mindestens drei Brettern mit einer
maximalen Dicke von $30\,\mathrm{mm}$. Die Bretter werden an den
Breitseiten verleimt; die Längsstöße werden keilverzinkt.


Ein Vorteil von Brettschichtholz besteht darin, dass Störungen wie Äste
oder Risse vor der Verarbeitung entfernt werden können. Dadurch sind
bessere Werte für die Tragfähigkeit möglich. 

Mit Brettschichtholz können Längen bis zu $35\,\mathrm{m}$ und Höhen
bis zu $2{,}2\,\mathrm{m}$ erreicht werden. Auch gekrümmte und
geknickte Träger sind herstellbar. 

\subsection{Massivholzplatten und Brettsperrholz}

Massivholzplatten bestehen aus Holzteilen mit einer Breite von mehr als
$18\,\mathrm{mm}$, die faserparallel miteinander verleimt werden.


Mehrschichtige Massivholzplatten bestehen in der Regel aus drei, vier
oder fünf Schichten. Der Aufbau ist symmetrisch und besitzt zwei
Decklagen sowie eine Mittellage. 

Brettsperrholz besteht aus einer ungeraden Anzahl von mindestens drei
Schichten. Jede zweite Schicht ist um $90^\circ$ gedreht. Dadurch sind
die Platten homogener und formstabiler. 

\subsection{Werkstoffe auf Furnierbasis}

Ein Furnier ist eine Platte mit einer Dicke von:

\begin{equation}
  0{,}5 \text{ bis } 8\,\mathrm{mm}
\end{equation}

Furniere werden durch Schälen oder Sägen gewonnen. Werden mehrere
Furnierlagen verleimt, entsteht Furnierschichtholz. 

Bei einem mindestens dreilagigen kreuzweise verleimten Aufbau wird von
Sperrholz gesprochen. Sperrholz besitzt eine höhere Formstabilität.


\subsection{Werkstoffe auf Spanbasis}

Holzwerkstoffe auf Spanbasis bestehen aus Holzspänen und anderen
holzartigen Faserstoffen. Durch Verleimen mit Kunstharz unter Wärme und
Druck entsteht ein plattenförmiger Werkstoff. 

Spanplatten können über ihren Querschnitt ein Dichteprofil besitzen, um
hohe Festigkeiten bei geringerem Gewicht zu ermöglichen. 

Aufgrund der verwendeten Kunstharze kann Formaldehyd an die
Umgebungsluft abgegeben werden. Für Bauholz und Möbelbau ist die
Formaldehydabgabe laut Dokument auf:

\begin{equation}
  0{,}1\,\mathrm{ppm}
\end{equation}

begrenzt. 

\subsection{Werkstoffe auf Faserbasis}

Holzfaserplatten entstehen durch Zusammenpressen und Entwässern eines
Holzfaserbreies. Das Holz wird zuvor unter Dampf mechanisch zerkleinert
und in Fasern aufgespalten. 

Die Einteilung erfolgt häufig nach Rohdichte:

\begin{itemize}
  \item Weichfaserplatten: $\rho \leq 450\,\mathrm{kg/m^3}$,
  \item mittelharte Holzfaserplatten,
  \item Hartfaserplatten: $\rho > 900\,\mathrm{kg/m^3}$.
\end{itemize}



\subsection{Verbundwerkstoffe}

Verbundwerkstoffe entstehen durch die Kombination verschiedener
Materialien. Dadurch können Vorteile einzelner Materialien kombiniert
und Schwächen ausgeglichen werden. 

Bei Holzverbundwerkstoffen bestehen die Baustoffe meist aus beidseitigen
Decklagen und einer Mittellage aus einem anderen Material.


\begin{longtable}{
  >{\raggedright\arraybackslash}p{0.30\textwidth}
  >{\raggedright\arraybackslash}p{0.50\textwidth}
}
\toprule
\textbf{Gruppe der Holzwerkstoffe} &
\textbf{Beispiele und Beschreibung} \\
\midrule
\endhead

Vollholzbasis &
Brettschichtholz, Massivholzplatten, Brettsperrholz. \\

Furnierbasis &
Furnierschichtholz, Sperrholz. \\

Spanbasis &
Spanplatten, OSB-Platten und weitere Spanwerkstoffe. \\

Faserbasis &
Weichfaserplatten, mittelharte Holzfaserplatten, Hartfaserplatten. \\

Verbundwerkstoffe &
Werkstoffe mit Decklagen und Mittellage aus unterschiedlichen
Materialien. \\

\bottomrule
\end{longtable}



\section{Beständigkeit}

Die Dauerhaftigkeit von Holz wird durch die Wahl der Holzart und durch
äußere Einflüsse bestimmt. 

Als sehr dauerhaft werden im Dokument Lärche, Eiche und Ulme genannt.
Fichte, Kiefer, Tanne und Buche besitzen eine gut ausgeprägte
Dauerhaftigkeit. Ahorn, Birke, Pappel und Weide schneiden hinsichtlich
der Dauerhaftigkeit schlecht ab. 

\begin{longtable}{
  >{\raggedright\arraybackslash}p{0.28\textwidth}
  >{\raggedright\arraybackslash}p{0.50\textwidth}
}
\toprule
\textbf{Dauerhaftigkeit} &
\textbf{Genannte Holzarten} \\
\midrule
\endhead

Sehr dauerhaft &
Lärche, Eiche, Ulme \\

Gut ausgeprägte Dauerhaftigkeit &
Fichte, Kiefer, Tanne, Buche \\

Geringe Dauerhaftigkeit &
Ahorn, Birke, Pappel, Weide \\

\bottomrule
\end{longtable}



Zu den holzartabhängigen Einflussfaktoren zählen:

\begin{itemize}
  \item Rohdichte,
  \item Härte,
  \item Kernholzanteil,
  \item Anteil widerstandsfähigeren Spätholzes,
  \item Porenanteil,
  \item Harzführung,
  \item Schnittfläche,
  \item Holzfeuchte,
  \item Oberflächenbehandlung,
  \item Holzschutzmittel.
\end{itemize}



Zu den äußeren Faktoren zählen Witterung, mechanische und chemische
Beanspruchungen, hohe Temperaturen und Brand sowie holzzerstörende
Organismen und Pilze. 

Die Beständigkeit gegenüber mechanischer Beanspruchung steigt mit
höherer Rohdichte, höherer Härte, größerem Kernholzanteil und
abnehmender Holzfeuchte. Außerdem besitzen stehende Jahresringe einen
Einfluss auf die Beständigkeit gegenüber mechanischer Beanspruchung.


\subsection{Chemische Beanspruchung}

Holz ist gegen chemische Beanspruchungen vergleichsweise beständig.
Bei pH-Werten unter $3$ oder über $10$ ist jedoch mit Holzzersetzung zu
rechnen. 

\subsection{Brandverhalten}

Holz ist ein brennbarer Baustoff. Die Abbrandgeschwindigkeit hängt von
Temperatur, Dichte, Störungen im Querschnitt wie Ästen und Rissen sowie
vom Feuchtegehalt ab. 

Bei der thermischen Zersetzung entsteht Holzkohle. Diese besitzt keine
relevanten Festigkeiten und ist nicht tragfähig. Die
Feuerwiderstandsfähigkeit hängt daher von der Eindringtiefe der
Holzkohle und der Resttragfähigkeit des verbleibenden Holzquerschnitts
ab. 

\begin{longtable}{
  >{\raggedright\arraybackslash}p{0.35\textwidth}
  >{\centering\arraybackslash}p{0.30\textwidth}
}
\toprule
\textbf{Werkstoff} &
\textbf{Abbrandgeschwindigkeit} \\
\midrule
\endhead

Eiche &
$0{,}50\,\mathrm{mm/min}$ \\

Fichte, auch brettschichtverleimt &
$0{,}65\,\mathrm{mm/min}$ \\

Kiefer &
$0{,}75\,\mathrm{mm/min}$ \\

Holzspannplatte &
$0{,}70$ bis $0{,}80\,\mathrm{mm/min}$ \\

\bottomrule
\end{longtable}



\subsection{Biologische Beanspruchung}

Biologische Holzeigenschaften beschreiben die Beständigkeit gegenüber
holzzerstörenden Organismen und Pilzen. Diese können Holz zersetzen und
die Eigenschaften des Bauholzes unerwünscht verändern.


Pilzbefall ist besonders bei einer Holzfeuchte von mehr als
$20\,\%$ und in einem Temperaturbereich von $3$ bis
$38\,^\circ\mathrm{C}$ möglich. 

Das Dokument unterscheidet holzfressende und holzbrütende Insekten.
Holzfressende Insekten ernähren sich vor allem im Larvenstadium von im
Bast enthaltenem Eiweiß. Holzbrütende Insekten leben meist in frischem
Holz und züchten Pilze als Nahrungsgrundlage. 

\section{Holzschutz}

Zur Erhöhung der Dauerhaftigkeit und Verbesserung der
Werkstoffeigenschaften werden Holzschutzmaßnahmen eingesetzt. Das
Dokument unterscheidet baulichen und chemischen Holzschutz.


\subsection{Baulicher Holzschutz}

Baulicher Holzschutz umfasst konstruktive Maßnahmen zur Erhöhung der
Lebensdauer oder Belastbarkeit. Dazu gehören insbesondere:

\begin{itemize}
  \item Auswahl geeigneter Hölzer,
  \item Dachüberstände,
  \item Sockelausbildungen,
  \item Tropfnasen.
\end{itemize}



Diese Maßnahmen sollen unter anderem verhindern, dass Wasser auf
Oberflächen oder in konstruktiven Bereichen stehen bleibt.

\subsection{Chemischer Holzschutz}

Chemischer Holzschutz zielt auf die Beständigkeit gegenüber
holzzerstörenden Organismen und Pilzen. Die Schutzmittel führen dem
Holz Stoffe zu, durch die es für Schädlinge ungenießbar wird.


\begin{longtable}{
  >{\raggedright\arraybackslash}p{0.30\textwidth}
  >{\raggedright\arraybackslash}p{0.50\textwidth}
}
\toprule
\textbf{Holzschutzart} &
\textbf{Ziel und Beispiele} \\
\midrule
\endhead

Baulicher Holzschutz &
Konstruktive Maßnahmen wie geeignete Holzwahl, Dachüberstände,
Sockelausbildungen und Tropfnasen. \\

Chemischer Holzschutz &
Verbesserung der Beständigkeit gegenüber holzzerstörenden Organismen
und Pilzen durch Schutzmittel. \\

\bottomrule
\end{longtable}



\section{Normengrundlage}

Für die Bemessung und Konstruktion von Holzbauten ist Eurocode~5
``Bemessung und Konstruktion von Holzbauten'' heranzuziehen.


Als österreichische Regelwerke nennt das Dokument:

\begin{itemize}
  \item ÖNORM EN 1995-1-1: Allgemeine Regeln und Regeln für den
  Hochbau,
  \item ÖNORM EN 1995-1-2: Allgemeine Regeln -- Bemessung für den
  Brandfall,
  \item ÖNORM EN 1995-2: Brücken.
\end{itemize}



\begin{merksatzbox}
\textbf{Merksatz:}
Für Holzbauten ist Eurocode~5 maßgebend. Die Festigkeitsklassen für
Bauholz werden nach EN~338 angegeben. Die Sortierung erfolgt visuell
oder maschinell.

\end{merksatzbox}
% =================================================
% KAPITEL 8: ASPHALT UND STRASSENBAU
% Diesen gesamten Block vor \end{document} einfügen.
% =================================================

\chapter{Asphalt und Straßenbau}

\section{Aufbau einer Straße}

Der Straßenaufbau wird in Oberbau, Unterbau und Untergrund gegliedert.
Der Oberbau umfasst die tragenden und deckenden Schichten der Straße.
Der Unterbau besteht aus Erdbau und bildet gemeinsam mit dem Untergrund
die Grundlage für den Straßenoberbau.


Bei Asphaltstraßen besteht der Oberbau aus mehreren Schichten. Von oben
nach unten sind dies üblicherweise:

\begin{itemize}
  \item Asphaltdeckschicht,
  \item Asphaltbinderschicht,
  \item Asphalttragschicht,
  \item ungebundene Tragschicht beziehungsweise Frostschutzschicht,
  \item Unterbau und Untergrund.
\end{itemize}



\begin{longtable}{
  >{\raggedright\arraybackslash}p{0.33\textwidth}
  >{\raggedright\arraybackslash}p{0.46\textwidth}
}
\toprule
\textbf{Schicht} &
\textbf{Aufgabe beziehungsweise Eigenschaft} \\
\midrule
\endhead

Asphaltdeckschicht &
Soll eben, griffig, dicht, verschleißfest und schubfest sein. \\

Asphaltbinderschicht &
Liegt zwischen Deckschicht und Asphalttragschicht. \\

Asphalttragschicht &
Tragende Asphaltschicht im Oberbau. \\

Ungebundene Tragschicht beziehungsweise Frostschutzschicht &
Tragfähigkeit und Frostsicherheit des Straßenaufbaus. \\

Unterbau und Untergrund &
Grundlage des gesamten Fahrbahnaufbaus. \\

\bottomrule
\end{longtable}



Die im Dokument dargestellten Schichtdicken betragen:

\begin{itemize}
  \item Asphaltdeckschicht: $1{,}5$ bis $6\,\mathrm{cm}$,
  \item Asphaltbinderschicht: $4$ bis $10\,\mathrm{cm}$,
  \item Asphalttragschicht: $8$ bis $20\,\mathrm{cm}$.
\end{itemize}



\begin{darstellungsbox}
\textbf{Darstellung: Fahrbahnaufbau einer Asphaltstraße}

Die Darstellung zeigt den Straßenaufbau von oben nach unten. An der
Oberfläche liegt die Asphaltdeckschicht. Darunter folgen
Asphaltbinderschicht und Asphalttragschicht. Unter den gebundenen
Asphaltschichten befinden sich ungebundene Tragschichten beziehungsweise
Frostschutzschichten. Den unteren Abschluss bilden Planum, Unterbau und
Untergrund.

Der Deckschicht werden die Eigenschaften eben, griffig, dicht,
verschleißfest und schubfest zugeordnet. Die Asphalttragschicht dient
vor allem der Tragfähigkeit.
\end{darstellungsbox}

\section{Unterbau und Bodenstabilisierung}

Der Unterbau gehört zum Erdbau. Böden können mit Bindemitteln wie Kalk,
Zement oder Mischbindemitteln stabilisiert werden. Diese Maßnahme wird
als Bodenverbesserung bezeichnet.


Durch Bodenstabilisierung können laut Dokument folgende Eigenschaften
verbessert werden:

\begin{itemize}
  \item Verdichtbarkeit,
  \item Tragfähigkeit,
  \item Setzungsverhalten.
\end{itemize}



Für das Unterbauplanum werden insbesondere folgende Anforderungen
angegeben:

\begin{equation}
  E_{v1} \geq 35\,\mathrm{MN/m^2}
\end{equation}

Der Wert $E_{v1}$ ist der Verformungsmodul bei Erstbelastung und wird
über einen Lastplattenversuch bestimmt.


Der Verdichtungsgrad $D_{\mathrm{Pr}}$ wird mit der Proctordichte
beschrieben:

\begin{equation}
  D_{\mathrm{Pr}}
  =
  \frac{\rho_a}{\rho_{\mathrm{Pr}}}
  \cdot 100
\end{equation}

Dabei bedeuten:

\begin{itemize}
  \item $D_{\mathrm{Pr}}$: erzielter Verdichtungsgrad in Prozent,
  \item $\rho_a$: erzielte Trockendichte,
  \item $\rho_{\mathrm{Pr}}$: Proctordichte.
\end{itemize}

Für den Unterbau wird folgende Anforderung genannt:

\begin{equation}
  D_{\mathrm{Pr}} > 100\,\%
\end{equation}



\section{Ungebundene und gebundene Tragschichten}

Ungebundene Tragschichten bestehen aus Gemischen natürlicher,
industriell hergestellter oder rezyklierter Gesteinskörnungen. Sie sind
nach RVS 08.15.01 geregelt.


Wesentliche Anforderungen an ungebundene Tragschichten sind:

\begin{itemize}
  \item Tragfähigkeit,
  \item Frostsicherheit,
  \item Verdichtbarkeit,
  \item ausreichender Verdichtungsgrad.
\end{itemize}



Die Tragfähigkeit wird über den Verformungsmodul $E_{v1}$ im
Lastplattenversuch beurteilt. Die Frostsicherheit wird über Korn- und
Mineralkriterien sowie durch Frosthebungsversuche beurteilt.


Gebundene Tragschichten werden eingesetzt, wenn Frostbeständigkeit oder
Tragfähigkeit ungebundener Tragschichten nicht ausreichend sind. Dabei
wird Gesteinskörnung mit einem Bindemittel gemischt.


Als Bindemittel nennt das Dokument:

\begin{itemize}
  \item Zement,
  \item Tragschichtbinder,
  \item Bitumen.
\end{itemize}



\section{Asphalt als Baustoff}

Asphalt besteht aus Gesteinskörnung und Bitumen. Das Dokument nennt für
Asphalt folgende Zusammensetzung:

\begin{itemize}
  \item $93$ bis $96\,\mathrm{M.\%}$ Gestein,
  \item $4$ bis $7\,\mathrm{M.\%}$ Bitumen,
  \item $0{,}5$ bis $15\,\mathrm{Vol.\%}$ Hohlraum.
\end{itemize}



Die Eigenschaften von Asphalt hängen insbesondere ab von:

\begin{itemize}
  \item den Eigenschaften der Gesteinskörnung,
  \item den Eigenschaften des Bitumens,
  \item der Verdichtung beim Einbau,
  \item der Temperatur beim Einbau,
  \item der Alterung des Asphalts.
\end{itemize}



\begin{merksatzbox}
\textbf{Merksatz:}
Asphalt besteht überwiegend aus Gesteinskörnung. Bitumen übernimmt die
Aufgabe des Bindemittels. Die Eigenschaften des fertigen Asphalts
werden wesentlich durch die Gesteinskörnung, das Bitumen sowie den
Einbau, insbesondere Verdichtung und Temperatur, beeinflusst.

\end{merksatzbox}

\section{Temperaturabhängiges Materialverhalten}

Das Materialverhalten von Asphalt ist temperaturabhängig. Bei tiefen
Temperaturen verhält sich Asphalt eher elastisch und ist rissanfällig.
Bei mittleren Temperaturen wird ein nichtlinear-elastisches Verhalten
mit möglicher Materialermüdung beschrieben.


Bei hohen Temperaturen verhält sich Asphalt viskos. Ab ungefähr
$50\,^\circ\mathrm{C}$ kann Asphalt teilweise plastisch werden. Dies
kann zur Bildung von Spurrinnen führen.


\begin{longtable}{
  >{\raggedright\arraybackslash}p{0.26\textwidth}
  >{\raggedright\arraybackslash}p{0.31\textwidth}
  >{\raggedright\arraybackslash}p{0.25\textwidth}
}
\toprule
\textbf{Temperaturbereich} &
\textbf{Materialverhalten} &
\textbf{Mögliche Folge} \\
\midrule
\endhead

Tiefe Temperaturen &
Eher elastisch &
Rissanfälligkeit und Rissbildung \\

Mittlere Temperaturen &
Nichtlinear elastisch &
Materialermüdung \\

Hohe Temperaturen &
Viskos; ab ungefähr $50\,^\circ\mathrm{C}$ teilweise plastisch &
Spurrinnenbildung \\

\bottomrule
\end{longtable}



Der Elastizitätsmodul von Asphalt ist ebenfalls temperaturabhängig. Für
Asphalt bei $18\,^\circ\mathrm{C}$ wird ein Wert von ungefähr

\begin{equation}
  E_{\mathrm{Asphalt}}
  \approx
  9000\,\mathrm{N/mm^2}
\end{equation}

angegeben.


Zum Vergleich nennt das Dokument ungefähr:

\begin{align}
  E_{\mathrm{Beton}}
  &\approx
  30000\,\mathrm{N/mm^2}, \\
  E_{\mathrm{Stahl}}
  &\approx
  210000\,\mathrm{N/mm^2}.
\end{align}



\section{Beanspruchung von Asphaltschichten}

Asphalt besitzt laut Dokument einen hohen Widerstand gegenüber
Druckbeanspruchung, versagt jedoch bei Zugbeanspruchung.


In der Asphaltdeckschicht ist vor allem Druckbeanspruchung relevant.
In der Asphalttragschicht ist insbesondere Zugbeanspruchung
maßgebend.


Der Verbund zwischen Deckschicht und Tragschicht ist wesentlich. Die
Schichten werden dafür verklebt.


\begin{longtable}{
  >{\raggedright\arraybackslash}p{0.32\textwidth}
  >{\raggedright\arraybackslash}p{0.47\textwidth}
}
\toprule
\textbf{Bereich} &
\textbf{Wesentliche Beanspruchung} \\
\midrule
\endhead

Asphaltdeckschicht &
Vor allem Druckbeanspruchung \\

Asphalttragschicht &
Vor allem Zugbeanspruchung \\

Verbund zwischen Deck- und Tragschicht &
Wesentlich für die Funktion des Aufbaus; wird durch Verklebung
hergestellt \\

\bottomrule
\end{longtable}



\section{Herstellung von Asphaltmischgut}

Asphaltmischgut wird in einer Asphaltmischanlage hergestellt. Der
Herstellungsprozess ist elektronisch gesteuert. Die einzelnen
Bestandteile des Asphaltmischguts werden dosiert, erwärmt und
miteinander vermischt.


Als zentrale Anlagenteile nennt das Dokument:

\begin{itemize}
  \item Doseure für die grobe Einwaage der Gesteinskörnung,
  \item Trockentrommel,
  \item Heißversiebung mit Mehrfachdecksieb,
  \item Mischturm,
  \item Mineralstoffwaage,
  \item Füllerwaage,
  \item Mischer,
  \item Bitumenlagertank,
  \item Bindemitteldosierung,
  \item Mischgutsilo.
\end{itemize}



In der Trockentrommel wird die Gesteinskörnung getrocknet und erwärmt.
Gleichzeitig wird ein Feinstaubanteil abgetrennt. Im Mischturm erfolgt
die genaue Dosierung der warmen Mineralstoffe; heißes Bitumen wird
eingedüst und mit den Mineralstoffen vermischt.


Die Mischzeit beträgt laut Dokument ungefähr:

\begin{equation}
  15\,\mathrm{s}
\end{equation}

Das frische Mischgut verlässt die Anlage meist mit einer Temperatur von:

\begin{equation}
  160 \text{ bis } 180\,^\circ\mathrm{C}
\end{equation}



\section{Asphaltmischgutarten}

Das Dokument unterscheidet Asphaltmischgutarten nach Einbau- und
Mischtemperatur.

\begin{longtable}{
  >{\raggedright\arraybackslash}p{0.27\textwidth}
  >{\centering\arraybackslash}p{0.25\textwidth}
  >{\raggedright\arraybackslash}p{0.27\textwidth}
}
\toprule
\textbf{Mischgutart} &
\textbf{Temperaturbereich} &
\textbf{Genannte Bindemittel} \\
\midrule
\endhead

Kaltmischgut und Kalteinbau &
Unter $30\,^\circ\mathrm{C}$ &
Bitumenemulsion, Fluxbitumen, polymermodifiziertes Bitumen \\

Warmeinbau &
$40$ bis $120\,^\circ\mathrm{C}$ &
Bitumenemulsion, Fluxbitumen, polymermodifiziertes Bitumen \\

Heißmischgut und Heißeinbau &
Über $130\,^\circ\mathrm{C}$ bis über $200\,^\circ\mathrm{C}$ &
Straßenbaubitumen, Fluxbitumen \\

\bottomrule
\end{longtable}



Der Heißeinbau ist laut Dokument das Standardverfahren. Kalt- und
Warmeinbau werden vor allem für Reparaturzwecke eingesetzt, etwa für
Belagssanierungen oder Spurrinnensanierungen mit Dünnschichtdecken.


\section{Mischkonzepte für Heißmischgut}

Für Heißmischgut werden unterschiedliche Mischkonzepte genannt.

\subsection{Mastix-Konzept}

Das Mastix-Konzept wird mit Gussasphalt verbunden. Es besitzt einen
hohen Bindemittelgehalt und praktisch keine Hohlräume. Die
Zuschlagsstoffe schwimmen in einer bituminösen Mörtelmatrix.


Weitere Merkmale sind:

\begin{itemize}
  \item kontinuierliche, feinkornangereicherte Kornverteilung,
  \item Lastabtrag hauptsächlich durch das Bindemittel,
  \item Einsatz vor allem bei Geh- und Fahrradwegen.
\end{itemize}



Für das Mastix-Konzept wird angegeben:

\begin{equation}
  B < 8{,}5\,\mathrm{M.\%}
\end{equation}

Der Hohlraumgehalt liegt bei ungefähr:

\begin{equation}
  H \approx 0\,\mathrm{Vol.\%}
\end{equation}



\subsection{Packungs-Konzept}

Das Packungs-Konzept wird auch als Betonprinzip bezeichnet und dem
Asphaltbeton zugeordnet. Es weist eine dichte Packung mit
kontinuierlicher Korngrößenverteilung auf.


Die Last wird vor allem über das Stützgerüst abgetragen. Gleichzeitig
soll ein möglichst dünner Bindemittelfilm ausgebildet werden.


Für das Packungs-Konzept wird angegeben:

\begin{align}
  B &< 6{,}5\,\mathrm{M.\%}, \\
  H &\approx 4\,\mathrm{Vol.\%}.
\end{align}



\subsection{Stützgerüst-Konzept}

Das Stützgerüst-Konzept wird auch Makadamprinzip genannt. Es beruht auf
einer Korn-zu-Korn-Abstützung der Gesteinskörner.


Wesentliche Merkmale sind:

\begin{itemize}
  \item diskontinuierliche Korngrößenverteilung,
  \item Verzahnung und innere Reibung des Korngerüsts,
  \item sekundäre stabilisierende Aufgabe des Bindemittels,
  \item Lastabtrag vor allem durch das Gesteinsgerüst.
\end{itemize}



Das Dokument unterscheidet zwei Arten:

\begin{itemize}
  \item dichtes Stützgerüst, beispielsweise Splittmastixasphalt,
  \item offenporiges Stützgerüst, beispielsweise Drainasphalt oder
  lärmdämmender Dünnschichtasphalt.
\end{itemize}



\begin{longtable}{
  >{\raggedright\arraybackslash}p{0.31\textwidth}
  >{\centering\arraybackslash}p{0.24\textwidth}
  >{\centering\arraybackslash}p{0.22\textwidth}
}
\toprule
\textbf{Mischkonzept} &
\textbf{Bindemittelgehalt} &
\textbf{Hohlraumgehalt} \\
\midrule
\endhead

Mastix-Konzept &
$B < 8{,}5\,\mathrm{M.\%}$ &
$H \approx 0\,\mathrm{Vol.\%}$ \\

Packungs-Konzept &
$B < 6{,}5\,\mathrm{M.\%}$ &
$H \approx 4\,\mathrm{Vol.\%}$ \\

Dichtes Stützgerüst &
$B < 7{,}5\,\mathrm{M.\%}$ &
$H \approx 4\,\mathrm{Vol.\%}$ \\

Offenporiges Stützgerüst &
$B < 5\,\mathrm{M.\%}$ &
$H \approx 20\,\mathrm{Vol.\%}$ \\

\bottomrule
\end{longtable}



\section{Anforderungen an Asphaltmischgut}

Das Dokument unterscheidet drei Gruppen von Anforderungen an
Asphaltmischgut:

\begin{itemize}
  \item allgemeine beziehungsweise volumetrische Anforderungen,
  \item konventionelle beziehungsweise empirische Anforderungen,
  \item gebrauchsverhaltensorientierte Anforderungen.
\end{itemize}



Allgemeine Anforderungen umfassen:

\begin{itemize}
  \item Rohdichte,
  \item Raumdichte,
  \item Hohlraum des Gesteinsgerüsts,
  \item Hohlraum im Asphalt,
  \item Bindemittelgehalt,
  \item Auffüllungsgrad.
\end{itemize}



Konventionelle Anforderungen beruhen auf Erfahrungswerten. Sie sind
einfach handhabbar, leiten jedoch keine Materialkenngrößen ab.


Gebrauchsverhaltensorientierte Anforderungen beruhen auf
realitätsnahen Beanspruchungen. Aus ihnen können Materialkenngrößen
abgeleitet werden, die Versuche sind jedoch komplexer.


\subsection{Konventionelle Prüfungen}

Als konventionelle Anforderungen nennt das Dokument:

\begin{itemize}
  \item Wasserempfindlichkeit nach EN 12697-12,
  \item Abrieb durch Spikereifen nach EN 12697-16,
  \item Affinität zwischen Gesteinskörnung und Bitumen nach
  EN 12697-11,
  \item Spurbildungstest nach EN 12697-22,
  \item Treibstoffbeständigkeit nach EN 12697-43.
\end{itemize}



\section{Marshall-Prüfung}

Die Marshall-Prüfung wird nach EN 12697-34 durchgeführt. Dazu werden
Marshall-Körper hergestellt.


Für die Herstellung wird loses Mischgut mit einer Temperatur von
$135\,^\circ\mathrm{C}$ in eine Stahlform gefüllt. Anschließend wird
der Probekörper im Marshall-Verdichtungsgerät auf jeder Seite mit
jeweils 50 Schlägen eines Fallhammers verdichtet.


Der Fallhammer fällt aus einer Höhe von:

\begin{equation}
  46\,\mathrm{cm}
\end{equation}



Mit der Marshall-Prüfung werden Marshall-Stabilität und Marshall-Fließwert
bestimmt.


Die Marshall-Stabilität ist der maximale Widerstand im
Kraft-Verformungs-Diagramm und wird in Kilonewton angegeben.


Der Marshall-Fließwert ist der Verformungsweg vom Beginn der Belastung
bis zum Erreichen der Höchstkraft. Er wird in Millimetern angegeben.


\begin{longtable}{
  >{\raggedright\arraybackslash}p{0.33\textwidth}
  >{\raggedright\arraybackslash}p{0.45\textwidth}
}
\toprule
\textbf{Kennwert} &
\textbf{Bedeutung} \\
\midrule
\endhead

Marshall-Stabilität &
Maximalwiderstand im Kraft-Verformungs-Diagramm; Einheit Kilonewton. \\

Marshall-Fließwert &
Verformungsweg vom Belastungsbeginn bis zum Erreichen der Höchstkraft;
Einheit Millimeter. \\

Marshall-Quotient &
Verhältnis von Marshall-Stabilität zu Marshall-Fließwert. \\

\bottomrule
\end{longtable}



\section{Asphalteinbau und Verdichtung}

Der Einbau von Asphalt erfolgt mit einem Fertiger. Der Fertiger leistet
bereits einen hohen Anteil der Verdichtungsarbeit. Die weitere
Nachverdichtung erfolgt mit Walzen.


Die Mindestverdichtung beträgt:

\begin{equation}
  98\,\%
  \text{ der Marshall-Dichte}
\end{equation}



Während des Einbauprozesses werden folgende Verdichtungsbereiche
angegeben:

\begin{longtable}{
  >{\raggedright\arraybackslash}p{0.43\textwidth}
  >{\centering\arraybackslash}p{0.29\textwidth}
}
\toprule
\textbf{Bereich beziehungsweise Arbeitsschritt} &
\textbf{Verdichtungsgrad} \\
\midrule
\endhead

Schneckenraum &
$74$ bis $78\,\%$ \\

Vorderwand beziehungsweise Stampfer &
$76$ bis $80\,\%$ \\

Glättblech &
$78$ bis $88\,\%$ \\

Bohle &
$86$ bis $90\,\%$ \\

Nachverdichtung mit Walzen &
$89$ bis $95\,\%$ \\

Weitere Nachverdichtung mit Walzen &
$92$ bis $98\,\%$ \\

\bottomrule
\end{longtable}



\section{Planung einer Asphaltbaustelle}

Asphalt muss heiß eingebaut werden. Die Temperaturgrenzen für das
Asphaltmischgut und insbesondere die unterste Einbautemperatur müssen
eingehalten werden. Für den Transport werden thermoisolierte
Transportfahrzeuge genannt.


Als Beispiele für Temperaturbereiche nennt das Dokument:

\begin{itemize}
  \item Bitumen \texttt{B70/100}: $140$ bis $180\,^\circ\mathrm{C}$,
  \item polymermodifiziertes Bitumen: $150$ bis $190\,^\circ\mathrm{C}$.
\end{itemize}



Für die Planung und Ausführung sind insbesondere folgende Punkte zu
berücksichtigen:

\begin{itemize}
  \item saubere und einwandfreie Unterlage,
  \item kontinuierlicher Einbau mit dem Fertiger,
  \item Sicherstellung des Schichtverbunds durch Vorspritzen,
  \item Längsnähte der Fertigerbahnen nicht in der Radspur,
  \item ausreichende Verdichtung,
  \item ausreichende Abkühlung vor der Verkehrsfreigabe.
\end{itemize}



Vor der Verkehrsfreigabe soll die Temperatur der Asphaltdecke unter

\begin{equation}
  35\,^\circ\mathrm{C}
\end{equation}

liegen. Das Dokument empfiehlt, die Asphaltdecke möglichst über Nacht
abkühlen zu lassen.


\begin{merksatzbox}
\textbf{Merksatz:}
Vor der Verkehrsfreigabe muss die Asphaltdecke ausreichend abkühlen.
Als Zielwert nennt das Dokument eine Asphalttemperatur von weniger als
$35\,^\circ\mathrm{C}$.

\end{merksatzbox}

\section{Qualitätsprüfungen und Erhaltung}

Während der Asphaltierungsarbeiten wird Mischgut entnommen. Nach der
Fertigstellung können zerstörende und zerstörungsfreie Verfahren zur
Qualitätsprüfung eingesetzt werden.


Zu den zerstörenden Verfahren zählt die Bohrkernentnahme. Diese erfolgt
laut Dokument ungefähr vier bis acht Wochen nach dem Einbau.


Bohrkerne dienen insbesondere zur Beurteilung von:

\begin{itemize}
  \item Schichtstärke,
  \item Schichtverbund,
  \item Haftzug,
  \item Mischguteigenschaften.
\end{itemize}



Eine Balkenentnahme wird nur in Sonderfällen durchgeführt.


Nach dem Bau folgt die Erhaltung. Regelmäßige Erhaltungsmaßnahmen sind
laut Dokument entscheidend für die Qualität einer Straße. Entscheidend
ist, mit der richtigen Maßnahme zum richtigen Zeitpunkt einzugreifen.


\begin{darstellungsbox}
\textbf{Darstellung: Zustandsverschlechterung und Erhaltung}

Die Darstellung zeigt die Verschlechterung des Straßenzustands mit
zunehmendem Alter. Ohne Erhaltungsmaßnahmen sinkt der Zustand
kontinuierlich.

Zwei Erhaltungsvarianten zeigen, dass rechtzeitig ausgeführte
Erhaltungsmaßnahmen den Zustand verbessern oder die weitere
Verschlechterung verzögern können. Als Ziel wird eine Steigerung der
Bauqualität und eine bessere Prognose des Straßenzustands genannt.
\end{darstellungsbox}

\section{Bituminöse Tragschichten im Bahnbau}

Im Bahnbau können bituminöse Tragschichten als wasserabweisende und
lastverteilende Konstruktionsschicht direkt unterhalb der
Schotterbettung eingesetzt werden.


Es handelt sich dabei um Heißmischgut mit folgenden Eigenschaften:

\begin{itemize}
  \item Schichtstärke: $9$ bis $12\,\mathrm{cm}$,
  \item Größtkorn: $32\,\mathrm{mm}$.
\end{itemize}



Die Anwendung wird insbesondere bei schlechten Bodenverhältnissen, aus
erschütterungstechnischen Gründen und auf Hochleistungsstrecken
genannt.


Als Vorteil wird eine lange Lebensdauer angeführt. Als Nachteil wird
genannt, dass der Unterbau anschließend nicht mehr mechanisch saniert
werden kann.


\section{Betonstraßenbau}

Neben Asphaltstraßen behandelt das Dokument Betonstraßen. Betonstraßen
weisen aufgrund ihrer mechanischen Eigenschaften eine andere
Druckverteilung als Asphaltkonstruktionen auf.


Für Betonstraßen wird ein Elastizitätsmodul von ungefähr

\begin{equation}
  E_{\mathrm{Beton}}
  \approx
  30000\,\mathrm{N/mm^2}
\end{equation}

angegeben. Der E-Modul wird als höher und konstanter als bei Asphalt
beschrieben.


In Österreich wird die unbewehrte Bauweise mit Betonplatten als
Regelbauweise genannt. Die Betondecke wird dabei durch Fugen in
einzelne Betonfelder unterteilt.


Querfugen erhalten Dübel zur Querkraftübertragung und zur Sicherung der
Höhenlage. Längsfugen werden mit Ankern zur Verbindung der Platten
ausgebildet.


Für die Betonfelder gelten laut Darstellung unter anderem folgende
Beziehungen:

\begin{align}
  L &\leq 1{,}5 \cdot B, \\
  L &\leq 25 \cdot D.
\end{align}

Dabei bedeuten:

\begin{itemize}
  \item $L$: Feldlänge,
  \item $B$: Feldbreite,
  \item $D$: Dicke der Betonplatte.
\end{itemize}



Zusätzlich wird eine maximale Feldlänge von:

\begin{equation}
  L \leq 5{,}5\,\mathrm{m}
\end{equation}

angegeben.


Bei extremen Temperaturen können sich Betonplatten ausdehnen und
verwerfen. Das Dokument bezeichnet dieses Phänomen als
\textit{Blow-up}. Die Oberseite einer Platte dehnt sich dabei stärker
aus als die Unterseite. Dadurch können Wölbspannungen entstehen.
Aufgrund des spröden Materialverhaltens kann es zu einem
explosionsartigen Aufplatzen kommen.


\begin{merksatzbox}
\textbf{Merksatz:}
Asphalt ist temperaturabhängig und kann bei hohen Temperaturen viskos
beziehungsweise teilweise plastisch werden. Betonstraßen bestehen
hingegen aus steiferen Betonplatten mit Fugen, Dübeln und Ankern.
\quelle{[3][^9]}
\end{merksatzbox}
% =================================================
% KAPITEL 9: MAUERWERK
% Diesen gesamten Block vor \end{document} einfügen.
% =================================================
% =================================================
% KAPITEL 9: MAUERWERK -- GRUNDLAGEN UND MAUERSTEINE
% Diesen gesamten Block vor \end{document} einfügen.
% =================================================

\chapter{Mauerwerk}

\section{Einleitung}

Mauerwerk ist ein Baukörper aus Mauersteinen und Mörtel, wobei der
Mörtel die Steine miteinander verbindet. Mauerwerk kann außerdem als
Verbundbaustoff verstanden werden, bei dem unterschiedliche
Mauersteine und Mörteltypen kombiniert werden, um für die jeweilige
Nutzung geeignete Eigenschaften zu erreichen.


Nach Eurocode~6 ist Mauerwerk definiert als:

\begin{quote}
\glqq Gefüge aus Mauersteinen, die in einem bestimmten Verband verlegt
und mit Mörtel verbunden worden sind.\grqq
\end{quote}



\begin{merksatzbox}
\textbf{Merksatz:}
Mauerwerk ist ein Verbundbaustoff aus Mauersteinen und Mörtel. Die
Mauersteine werden in einem bestimmten Verband verlegt und durch Mörtel
miteinander verbunden.

\end{merksatzbox}

\section{Aufgaben von Mauerwerksbauteilen}

Mauerwerksbauteile übernehmen statische, bauphysikalische und
architektonische Aufgaben.

\begin{longtable}{
  >{\raggedright\arraybackslash}p{0.29\textwidth}
  >{\raggedright\arraybackslash}p{0.50\textwidth}
}
\toprule
\textbf{Aufgabenbereich} &
\textbf{Aufgaben von Mauerwerk} \\
\midrule
\endhead

Statische Aufgaben &
Ableiten von Lasten, Ausbildung von End- oder Zwischenauflagern,
Aussteifung des gesamten Bauwerks sowie rissfreie Aufnahme von
Verformungen. \\

Bauphysikalische Aufgaben &
Wärmeschutz, Schallschutz, Brandschutz und Witterungsschutz. \\

Architektonische Aufgaben &
Gestaltung von Oberflächen, optische Gliederung von Bauwerken und
Raumabschluss. \\

\bottomrule
\end{longtable}



\section{Mauersteine}

Ein Mauerstein ist ein vorgefertigtes Bauteil zur Verwendung im
Mauerwerksbau. Mauersteine können grundsätzlich in natürliche und
künstlich hergestellte Steine eingeteilt werden.


Natursteine werden durch Bearbeitung einer Fachkraft für den
Mauerwerksbau vorbereitet. Bei künstlich hergestellten Steinen erfolgt
die Herstellung in einem industriellen Herstellungsprozess.


\subsection{Natürliche Steine}

Natursteinmauerwerk besteht aus natürlichen Steinen beziehungsweise
Naturbruchsteinen. Natursteine können bearbeitet oder weitgehend
unbearbeitet verwendet werden. Die Wahl der Steinart hängt vom
Vorkommen, vom Einsatzgebiet und vom gewünschten Mauerwerkstyp ab.


Das Dokument nennt unter anderem folgende Arten von
Natursteinmauerwerk:

\begin{itemize}
  \item Trockenmauerwerk,
  \item Zyklopenmauerwerk,
  \item Bruchsteinmauerwerk,
  \item unregelmäßiges Schichtenmauerwerk,
  \item hammerrechtes Schichtenmauerwerk,
  \item regelmäßiges Schichtenmauerwerk.
\end{itemize}



Trockenmauerwerk wird ohne Mörtel errichtet. Es besteht in vielen Fällen
aus Natursteinen oder Naturbruchsteinen. Heute wird es vor allem in
Gartenanlagen, bei groben steinigen Hangsicherungen und im Wasserbau
eingesetzt.


\subsection{Künstlich hergestellte Steine}

Künstlich hergestellte Mauersteine werden in keramische und mineralisch
gebundene Steine unterteilt.

\begin{longtable}{
  >{\raggedright\arraybackslash}p{0.32\textwidth}
  >{\raggedright\arraybackslash}p{0.47\textwidth}
}
\toprule
\textbf{Gruppe} &
\textbf{Beispiele} \\
\midrule
\endhead

Keramische Steine &
Mauerziegel, Klinker und feuerfeste keramische Baustoffe. \\

Mineralisch gebundene Steine &
Kalksandsteine, Hüttensteine, Gasbetonsteine und Betonsteine. \\

\bottomrule
\end{longtable}



\section{Eigenschaften von Mauersteinen}

Die Eigenschaften von Mauersteinen hängen vom Zusammenspiel der
Rohstoffe, der Materialstruktur und des Herstellungsverfahrens ab.
Diese Faktoren beeinflussen Gewicht, Rohdichte, statische und
bauphysikalische Eigenschaften sowie die Eignung für bestimmte
Anwendungsbereiche.


\subsection{Rohdichte}

Die Rohdichte beschreibt das Raumgewicht eines Mauersteins.
Hohlkammern, Löcher, Nuten, Federn und Schlitze werden dabei
berücksichtigt.


Die Rohdichte wird in \(\mathrm{kg/dm^3}\) angegeben. Die
Rohdichteklassen reichen bei genormten Steinen von:

\begin{equation}
  0{,}5 \text{ bis } 2{,}2\,\mathrm{kg/dm^3}
\end{equation}

Die Zahlenwerte beschreiben jeweils die obere Grenze der
Rohdichteklasse.


Die Rohdichte hat einen wesentlichen Einfluss auf:

\begin{itemize}
  \item die Wärmedämmung,
  \item die Schalldämmung,
  \item das Stückgewicht,
  \item die Arbeitsleistung bei der Errichtung eines Bauwerks.
\end{itemize}



\subsection{Druckfestigkeit}

Die Festigkeit ist bei Mauersteinen eine wesentliche Kenngröße. Sie
beeinflusst gemeinsam mit der Festigkeit des Mörtels die
Gesamtfestigkeit der Wände eines Bauwerks.


Die Festigkeit wird als Druckfestigkeit in \(\mathrm{N/mm^2}\)
beziehungsweise \(\mathrm{MPa}\) angegeben. Sie wird durch einen
Druckversuch bestimmt.


Bei genormten Ziegeln ist die Druckfestigkeit in Klassen eingeteilt.
Die Klassenbezeichnung entspricht dabei dem geringsten Wert, den ein
Ziegel dieser Klasse aufweisen muss.


\begin{longtable}{
  >{\centering\arraybackslash}p{0.25\textwidth}
  >{\centering\arraybackslash}p{0.25\textwidth}
  >{\centering\arraybackslash}p{0.25\textwidth}
}
\toprule
\textbf{Druckfestigkeitsklasse} &
\textbf{Mittelwert} &
\textbf{Kleinster Einzelwert} \\
\midrule
\endhead

2 &
$2{,}5\,\mathrm{N/mm^2}$ &
$2{,}0\,\mathrm{N/mm^2}$ \\

4 &
$5{,}0\,\mathrm{N/mm^2}$ &
$4{,}0\,\mathrm{N/mm^2}$ \\

6 &
$7{,}5\,\mathrm{N/mm^2}$ &
$6{,}0\,\mathrm{N/mm^2}$ \\

12 &
$15{,}0\,\mathrm{N/mm^2}$ &
$12{,}0\,\mathrm{N/mm^2}$ \\

20 &
$25{,}0\,\mathrm{N/mm^2}$ &
$20{,}0\,\mathrm{N/mm^2}$ \\

28 &
$35{,}0\,\mathrm{N/mm^2}$ &
$28{,}0\,\mathrm{N/mm^2}$ \\

36 &
$45{,}0\,\mathrm{N/mm^2}$ &
$36{,}0\,\mathrm{N/mm^2}$ \\

48 &
$60{,}0\,\mathrm{N/mm^2}$ &
$48{,}0\,\mathrm{N/mm^2}$ \\

60 &
$75{,}0\,\mathrm{N/mm^2}$ &
$60{,}0\,\mathrm{N/mm^2}$ \\

\bottomrule
\end{longtable}



\subsection{Stückgewicht}

Das Stückgewicht eines Mauersteins wird wesentlich durch Art, Anzahl
und Größe der Hohlräume beeinflusst. Der Hohlraumanteil beeinflusst
damit die Rohdichte und das Gewicht eines einzelnen Steins.


Ein geringes Stückgewicht bei gleichzeitig großem Volumen kann die
ständigen Lasten aus Eigengewicht reduzieren und die Arbeitsleistung bei
der Herstellung erhöhen.


\subsection{Frostbeständigkeit}

Die Frostbeständigkeit beeinflusst wesentlich das Einsatzgebiet eines
Mauersteins. Bereits bei der Planung müssen die Exposition und ein
möglicher Frostangriff berücksichtigt werden, damit spätere Schäden am
Bauwerk vermieden werden.


Ob ein Stein frostbeständig ist, muss durch eine eigene Prüfung
nachgewiesen werden. Frostbeständige Steine sind beispielsweise
Klinker, Vormauerziegel sowie bestimmte Kalksand- und Hüttensteine.


\begin{merksatzbox}
\textbf{Merksatz:}
Rohdichte, Druckfestigkeit, Stückgewicht und Frostbeständigkeit sind
zentrale Eigenschaften von Mauersteinen. Sie beeinflussen sowohl die
bauphysikalische Eignung als auch die Tragfähigkeit und Verarbeitung.

\end{merksatzbox}

\section{Mauersteinformate und Steinformen}

Die häufig verwendeten Steinformate mit festen Abmessungen sind das
Dünnformat \texttt{DF} und das Normalformat \texttt{NF}.

\begin{longtable}{
  >{\raggedright\arraybackslash}p{0.30\textwidth}
  >{\centering\arraybackslash}p{0.42\textwidth}
}
\toprule
\textbf{Format} &
\textbf{Nennmaß: Länge $\times$ Breite $\times$ Höhe} \\
\midrule
\endhead

Dünnformat \texttt{DF} &
$240 \times 115 \times 52\,\mathrm{mm}$ \\

Normalformat \texttt{NF} &
$240 \times 115 \times 71\,\mathrm{mm}$ \\

\bottomrule
\end{longtable}



Mauersteine können verschiedene Formen besitzen. Diese ergeben sich aus
unterschiedlichen Einsatzgebieten und technologischen Entwicklungen.

\begin{itemize}
  \item Vollstein,
  \item Vollstein mit Griffschlitz,
  \item Blockstein,
  \item Hohlblockstein,
  \item Langlochziegel,
  \item Hochlochstein mit versetzten Stegen,
  \item Hochlochstein mit durchlaufenden Stegen.
\end{itemize}



\section{Keramisch gebundene Steine: Ziegel}

Ziegel können allgemein als \glqq gebrannte Erden\grqq{} bezeichnet
werden. Der Brennvorgang bildet die gewünschten Eigenschaften aus oder
verstärkt sie.


Als Ausgangsmaterial dient Ton, vor allem jedoch Lehm. Lehm ist eine
Mischung aus Sand, Ton und Schluff. Diese Bestandteile entstehen durch
die Verwitterung von Gestein.


Durch die Einlagerung von Wasser in die schichtig aufgebauten
Tonminerale ist Lehm im feuchten Zustand gut formbar. Diese Eigenschaft
ermöglicht die Herstellung unterschiedlicher Ziegelformen und Lochungen.


Die Wasseraufnahme bewirkt zugleich, dass beim Brennen ein
Schrumpfprozess einsetzt. Nicht bildsame Bestandteile wie Sand und Kalk
wirken sich günstig auf das Trocknungs- und Schwindverhalten aus.


\subsection{Porosierungsstoffe}

Neben den natürlichen Rohstoffen können Porosierungsstoffe zugesetzt
werden. Das Dokument nennt beispielsweise Sägespäne oder Schalen von
Sonnenblumenkernen.


Diese Stoffe verglühen beim Brennen und hinterlassen kleine
luftgefüllte Hohlräume. Die erhöhte Porosität beeinflusst die
bauphysikalischen Eigenschaften des fertigen Ziegels wesentlich.


\subsection{Herstellung von Ziegeln}

Der Brennvorgang ist der letzte Schritt der Ziegelherstellung. Dabei
verliert der Ziegel zunächst sein gebundenes Wasser, sowohl
physikalisch als auch chemisch gebundenes Wasser.


Brenntemperatur und Gleichmäßigkeit des Brandes haben einen
wesentlichen Einfluss auf die Eigenschaften des Ziegels.


Ab ungefähr folgender Temperatur kommt es zur Neubildung von
Kristallen:

\begin{equation}
  800\,^\circ\mathrm{C}
\end{equation}

Die neu gebildeten Kristalle vernetzen sich und geben dem Ziegel
Festigkeit und Zusammenhalt.


Bei Temperaturen von ungefähr

\begin{equation}
  1000 \text{ bis } 1300\,^\circ\mathrm{C}
\end{equation}

weichen sich Tonminerale an der Ziegeloberfläche auf und verschmelzen
miteinander. Dadurch wird die Wasseraufnahmefähigkeit vermindert.
Dieser Vorgang wird Sintern beziehungsweise Dichtbrennen genannt.


Feuerfeste Erzeugnisse werden bei Temperaturen zwischen

\begin{equation}
  1300 \text{ und } 1800\,^\circ\mathrm{C}
\end{equation}

gebrannt. Sie werden beispielsweise für Kamine und Ofenverkleidungen
verwendet.


\begin{darstellungsbox}
\textbf{Darstellung: Brennen von Ziegeln}

Beim Brennen verliert der Ziegel zunächst physikalisch und chemisch
gebundenes Wasser. Ab etwa $800\,^\circ\mathrm{C}$ bilden sich neue
Kristalle, die dem Ziegel Festigkeit und Zusammenhalt verleihen.

Bei etwa $1000$ bis $1300\,^\circ\mathrm{C}$ können die
Tonminerale an der Oberfläche aufweichen und miteinander
verschmelzen. Dieser Prozess wird als Sintern oder Dichtbrennen
bezeichnet und vermindert die Wasseraufnahmefähigkeit.

Feuerfeste keramische Erzeugnisse werden bei Temperaturen von etwa
$1300$ bis $1800\,^\circ\mathrm{C}$ hergestellt.
\end{darstellungsbox}

\subsection{Tunnelofen}

Früher wurden Ringöfen verwendet. Dabei hing die Ziegelqualität stark
vom Ort des Ziegels innerhalb des Ofens ab, weil Brenntemperatur und
Dauer der Einwirkung der Maximaltemperatur unterschiedlich waren.


Heute werden Tunnelöfen verwendet. Die Ziegel werden auf einem
Förderband durch den Ofen transportiert. Die Befeuerung liegt in der
Mitte des Tunnels. Dadurch erfolgen Temperaturanstieg und
Temperaturabfall langsam und schrittweise.


Durch den kontrollierten Temperaturverlauf werden Zwangsspannungen und
die Gefahr von Spannungsrissen reduziert.


\section{Mineralisch gebundene Steine}

Mineralisch gebundene Steine unterscheiden sich von Ziegeln dadurch,
dass nicht die gesamte Masse einem Brennprozess unterzogen wird.
Üblicherweise wird das Bindemittel als gebranntes Erzeugnis zugesetzt.
In manchen Fällen kann auch die Körnung gebrannt sein.


Zu den mineralisch gebundenen Steinen zählen unter anderem
Kalksandsteine, Porenbetonsteine, Hüttensteine und Betonsteine.


\subsection{Kalksandsteine}

Kalksandsteine bestehen aus Kalk, kieselsäurehaltigem Sand und Wasser.
Das Wasser dient zum Ablöschen des gebrannten Kalks.


Das Gemisch wird unter Druck zu Rohlingen geformt. Anschließend müssen
die Rohlinge gehärtet werden. Die Härtung erfolgt in Autoklaven bei
einem Dampfdruck von ungefähr:

\begin{equation}
  16 \text{ bis } 21\,\mathrm{bar}
\end{equation}

und einer Temperatur zwischen:

\begin{equation}
  160 \text{ und } 250\,^\circ\mathrm{C}
\end{equation}



Der hochgespannte Dampf schließt Kieselsäure an der Oberfläche der
Sandkörner chemisch auf. Dadurch verbindet sich Kalk mit den
Zuschlägen. Die entstehenden Kalkhydrosilikate sind wasserbeständig und
weisen hohe Festigkeiten auf.


\subsection{Porenbeton}

Porenbeton wird auch Gasbeton genannt. Seine Hauptbestandteile sind
feingemahlener Quarzsand, Bindemittel wie Zement oder Kalk, Wasser sowie
ein Treibmittel. Als Treibmittel wird meist Aluminiumpulver verwendet.


Zusätzlich können Gips oder Anhydrit zugesetzt werden, um den
Produktionsablauf zu steuern.


Das flüssige Gemisch wird in Formen gegossen. Das Treibmittel reagiert
mit Calciumhydroxid und bildet Wasserstoffgas. Dadurch geht die Masse
auf und es entsteht eine poröse Struktur.


Nach dem Aufgehen kann der Porenbeton auf die gewünschte Größe
geschnitten und anschließend im Autoklav dampfgehärtet werden.


\section{Mauermörtel}

Mörtel ist eine Mischung aus anorganischen Bindemitteln, Zuschlägen und
Wasser. Zusatzstoffe und Zusatzmittel können gegebenenfalls enthalten
sein.


Mauermörtel verbindet Mauersteine kraftschlüssig zu einem Mauerwerk. Er
ermöglicht dadurch die Aufnahme und Ableitung von Druck-, Zug- und
Momentenbeanspruchungen.


Zusätzlich dient Mauermörtel als vollflächiges Lager zwischen den
Mauersteinen. Dadurch trägt er zu Feuchtigkeits-, Wärme-, Schall- und
Brandschutz bei.


Mörtel gleicht außerdem Maßabweichungen der Mauersteine aus. Je geringer
die Fugendicke ist, desto höher sind die Anforderungen an die
Maßgenauigkeit der Steine.


\begin{longtable}{
  >{\raggedright\arraybackslash}p{0.34\textwidth}
  >{\raggedright\arraybackslash}p{0.46\textwidth}
}
\toprule
\textbf{Mörtelart nach Bindemittel} &
\textbf{Erhärtungsverhalten} \\
\midrule
\endhead

Kalkmörtel &
Nicht hydraulisch. \\

Lehmmörtel &
Nicht hydraulisch. \\

Gipsmörtel &
Nicht hydraulisch. \\

Zementmörtel &
Hydraulisch. \\

Kalkzementmörtel &
Hydraulisch. \\

\bottomrule
\end{longtable}



Hydraulische Mörtel benötigen das Wasser aus der Mörtelmischung zum
Erhärten. Nicht hydraulische Mörtel trocknen und erhärten an der Luft.


\subsection{Fugendicken}

Mauermörtel wird üblicherweise in einer Schichtdicke von etwa
\(1\) bis \(2\,\mathrm{cm}\) aufgetragen.


Dünnbettmörtel besitzt eine maximale Schichtdicke von:

\begin{equation}
  3\,\mathrm{mm}
\end{equation}

Er wird daher auch als Klebemörtel bezeichnet.


Dünnbettmörtel wird insbesondere bei Gasbeton, Kalksandstein und plan
geschliffenen Steinen verwendet. Dies setzt geringe Abweichungen der
Steinabmessungen voraus.


\section{Klassifizierung von Mauerwerk}

Mauerwerk kann nach verschiedenen Gesichtspunkten klassifiziert werden.

\begin{longtable}{
  >{\raggedright\arraybackslash}p{0.30\textwidth}
  >{\raggedright\arraybackslash}p{0.50\textwidth}
}
\toprule
\textbf{Kriterium} &
\textbf{Einteilung} \\
\midrule
\endhead

Steine &
Natursteinmauerwerk oder Ziegelmauerwerk. \\

Statische Funktion &
Tragend, nicht tragend, bewehrt, nicht bewehrt, aussteifende Wand oder
Schubwand. \\

Zusammensetzung &
Trockenmauerwerk, Mörtelmauerwerk, homogenes oder inhomogenes
Mauerwerk. \\

Sichtbarkeit &
Sichtmauerwerk, Verblendmauerwerk, verputztes oder wärmegedämmtes
Mauerwerk. \\

Art &
Einschalige Wand, zweischalige Wand oder Verbundwand. \\

\bottomrule
\end{longtable}



\subsection{Tragendes und nicht tragendes Mauerwerk}

Eine tragende Wand ist dafür vorgesehen, neben dem Eigengewicht auch
Lasten aus anderen Bauteilen schadlos in den Untergrund oder in andere
tragende Bauteile abzuleiten.


Nicht tragende Mauerwerksbauteile müssen nur Eigenlasten und direkt auf
sie aufgebrachte Lasten aufnehmen. Sie dienen beispielsweise als
Raumtrennung, Ausfachung oder Verblendmauerwerk.


\subsection{Einschalige, zweischalige und Verbundwände}

Bei Verbundwänden ergänzt eine Dämmschicht die einschalige Wand. Die
Wand übernimmt dabei Lastabtragung, Schallschutz und Wärmespeicherung,
während die Dämmschicht die Dämmeigenschaften verbessert.


Bei zweischaligen Wänden sind die Funktionen getrennt:

\begin{itemize}
  \item Die Innenschale ist das tragende Element.
  \item Die Zwischenschicht übernimmt die Dämmfunktion.
  \item Die Außenschale übernimmt Witterungs- und Frostschutz.
\end{itemize}



\section{Druckfestigkeit von Mauerwerk}

Die Tragfähigkeit von Mauerwerk unter Druckbeanspruchung ist im
Vergleich zu anderen Beanspruchungsarten am höchsten. Mauerwerk wird
daher vor allem für druckbeanspruchte Bauteile verwendet.


Mauerwerk ist ein zusammengesetzter Baustoff aus mehreren Materialien
mit unterschiedlichen Eigenschaften. Das Bruchbild ist daher komplex.
Ein Bruch erfolgt im Allgemeinen durch Versagen der Mauersteine infolge
vertikaler Spaltzugrisse.


Die tatsächlich ableitbare Druckspannung wird unter anderem beeinflusst
durch:

\begin{itemize}
  \item Feuchtezustand,
  \item Lagerfugendicke,
  \item Lochbild des Steins,
  \item Ausführungsqualität.
\end{itemize}



Die Druckfestigkeit des Mauerwerks kann vereinfacht mit folgender
Gleichung ermittelt werden:

\begin{equation}
  f_k =
  K \cdot f_b^a \cdot f_m^b
\end{equation}

Dabei bedeuten:

\begin{itemize}
  \item \(f_k\): Druckfestigkeit des Mauerwerksverbands in
  \(\mathrm{MPa}\),
  \item \(f_b\): Steindruckfestigkeit in \(\mathrm{MPa}\),
  \item \(f_m\): Mörteldruckfestigkeit in \(\mathrm{MPa}\),
  \item \(K\), \(a\), \(b\): Beiwerte abhängig von Mauerstein und
  Mörtel.
\end{itemize}



\section{Schubfestigkeit von Mauerwerk}

Bei Schub- und Scherbeanspruchung kann zwischen Scheibenbeanspruchung
und Plattenbeanspruchung unterschieden werden.


Die Art des Versagens hängt von der vorhandenen Druckspannung im
Mauerwerk ab:

\begin{itemize}
  \item Bei geringer Druckspannung können Lager- und Stoßfugen
  versagen.
  \item Bei hoher Druckspannung kann ein kombiniertes Versagen von
  Stoßfugen und Steinen auftreten.
  \item Bei sehr hoher Druckspannung kann reines Steinversagen
  auftreten.
\end{itemize}



Die Schubfestigkeit von Mauerwerk wird durch die Festigkeitseigenschaften
der Mauersteine, die Verbundfestigkeit und die Ausführungsqualität
beeinflusst.


Für vollständig vermörtelte Fugen lautet die Berechnung der
Schubfestigkeit:

\begin{equation}
  f_{vk} =
  f_{vko} + 0{,}4 \cdot \sigma_d
  \leq
  0{,}065 \cdot f_b
\end{equation}

Für knirsch gestoßene Stirnflächen lautet sie:

\begin{equation}
  f_{vk} =
  0{,}5 \cdot f_{vko} + 0{,}4 \cdot \sigma_d
  \leq
  0{,}045 \cdot f_b
\end{equation}

Dabei bedeuten:

\begin{itemize}
  \item \(f_{vk}\): Schubfestigkeit in \(\mathrm{MPa}\),
  \item \(f_{vko}\): Anfangsscherfestigkeit, abhängig vom Mörtel,
  \item \(\sigma_d\): Druckspannung in \(\mathrm{MPa}\),
  \item \(f_b\): Steindruckfestigkeit.
\end{itemize}



\section{Elastizitätsmodul von Mauerwerk}

Der Elastizitätsmodul beschreibt die Beziehung zwischen Spannung und
Dehnung eines Materials. Für Mauerwerk wird der E-Modul nach der
folgenden Beziehung berechnet:

\begin{equation}
  E =
  1000 \cdot f_k
\end{equation}

Dabei ist \(f_k\) die Druckfestigkeit des Mauerwerks.


\begin{merksatzbox}
\textbf{Merksatz:}
Die Mauerwerksdruckfestigkeit hängt von der Steindruckfestigkeit, der
Mörteldruckfestigkeit und von material- und mörtelabhängigen
Beiwerten ab. Die Schubfestigkeit wird zusätzlich durch die
Druckspannung, die Verbundfestigkeit und die Ausführungsqualität
beeinflusst.
\quelle{[16][17]}
\end{merksatzbox}
% =================================================
% KAPITEL 10: NEUE, BIOBASIERTE UND
% KREISLAUFFÄHIGE MATERIALIEN
% Diesen gesamten Block vor \end{document} einfügen.
% =================================================

\chapter{Neue, biobasierte und kreislauffähige Materialien}

\section{Kreislauffähige Baustoffe und Rezyklate}

Das Dokument behandelt Rezyklate als ``Baustoff der Zukunft''. Dabei
werden Beispiele für die Materialaufbereitung angeführt.


Die im Dokument dargestellten Beispiele beziehen sich auf die
Materialaufbereitung durch Hengl in Maissau sowie auf eine
Materialaufbereitung in Wien.


Rezyklate stehen im Zusammenhang mit der Wiederverwendung und
Aufbereitung von Baustoffen. Im Straßenbau werden beispielsweise
rezyklierte Gesteinskörnungen für ungebundene Tragschichten genannt.
Ungebundene Tragschichten können aus natürlichen, industriell
hergestellten und rezyklierten Gesteinskörnungen bestehen.


\begin{merksatzbox}
\textbf{Merksatz:}
Rezyklate werden im Dokument als Baustoff der Zukunft bezeichnet.
Rezyklierte Gesteinskörnungen können unter anderem für ungebundene
Tragschichten im Straßenbau verwendet werden.

\end{merksatzbox}

\section{Rezyklierte Gesteinskörnungen im Straßenbau}

Ungebundene Tragschichten bestehen aus Gemischen von
Gesteinskörnungen. Das Dokument nennt dafür natürliche, industriell
hergestellte und rezyklierte Gesteinskörnungen.


Die wesentlichen Anforderungen an ungebundene Tragschichten betreffen
insbesondere:

\begin{itemize}
  \item Tragfähigkeit,
  \item Frostsicherheit,
  \item Verdichtbarkeit,
  \item Verdichtungsgrad.
\end{itemize}



Die Tragfähigkeit wird über den Verformungsmodul $E_{v1}$ im
Lastplattenversuch beurteilt. Die Frostsicherheit wird über
Korngrößen- und Mineralkriterien sowie über Frosthebungsversuche
beurteilt.


\begin{darstellungsbox}
\textbf{Darstellung: Rezyklate als Baustoff}

Die Darstellung zeigt die Materialaufbereitung von Rezyklaten. Der
Begriff Rezyklate wird dabei im Dokument als Teil des Themenbereichs
``Baustoff der Zukunft'' verwendet.

Für den Straßenbau nennt das Dokument rezyklierte Gesteinskörnungen als
möglichen Bestandteil ungebundener Tragschichten. Dabei müssen die
Anforderungen an Tragfähigkeit, Frostsicherheit, Verdichtbarkeit und
Verdichtungsgrad erfüllt werden.
\end{darstellungsbox}

\section{Stampflehm}

Stampflehm wird im Dokument als Lehmbaustoff behandelt. Die
Zusammensetzung von Lehm wird mit folgenden Bestandteilen angegeben:

\begin{itemize}
  \item Ton,
  \item Schluff,
  \item Sand,
  \item Kies.
\end{itemize}



Für Stampflehm verweist das Dokument auf Erfahrungswerte gemäß dem
``Design Guide für Stampflehm'' von Erden.


\begin{longtable}{
  >{\raggedright\arraybackslash}p{0.31\textwidth}
  >{\raggedright\arraybackslash}p{0.49\textwidth}
}
\toprule
\textbf{Bestandteil} &
\textbf{Zuordnung} \\
\midrule
\endhead

Ton &
Bestandteil von Lehm \\

Schluff &
Bestandteil von Lehm \\

Sand &
Bestandteil von Lehm \\

Kies &
Bestandteil von Lehm \\

\bottomrule
\end{longtable}



\begin{merksatzbox}
\textbf{Merksatz:}
Stampflehm besteht aus Lehm. Als Bestandteile von Lehm nennt das
Dokument Ton, Schluff, Sand und Kies.

\end{merksatzbox}

\section{Zusammenfassung}

Die vorliegenden Dokumentauszüge behandeln Rezyklate als Baustoff der
Zukunft und zeigen deren Materialaufbereitung. Im Straßenbau können
rezyklierte Gesteinskörnungen in ungebundenen Tragschichten eingesetzt
werden. \quelle{[1][^3]}

Als weiterer Werkstoff wird Stampflehm genannt. Dieser basiert auf Lehm
mit den Bestandteilen Ton, Schluff, Sand und Kies.


\begin{merksatzbox}
\textbf{Merksatz:}
Kreislauffähige Baustoffe umfassen im Dokument insbesondere Rezyklate.
Als mineralischer Lehmbaustoff wird Stampflehm mit den Bestandteilen
Ton, Schluff, Sand und Kies behandelt.
\quelle{[1][^2]}
\end{merksatzbox}
\end{document}